\documentclass[11pt]{amsart}
\usepackage[T1]{fontenc}
\usepackage{lmodern}
\input{glyphtounicode}
\pdfgentounicode=1
\usepackage{amsmath,amssymb}
\usepackage[table]{xcolor}
\usepackage{array}
\usepackage{booktabs}
\usepackage[margin=1in]{geometry}
\usepackage{url}
\usepackage{listings}
\usepackage[pdfauthor={311859884+mt0-svg@users.noreply.github.com},
            pdftitle={The plane quartic of Furio and Lombardo has exactly four rational points},
            pdfkeywords={rational points, plane quartic, twist of the Klein quartic, Selmer group Chabauty, Prym variety, genus 2 Jacobian, 2-adic logarithm, 7-adic Galois representations, Lean 4, formal verification},
            bookmarksnumbered=true]{hyperref}
\hypersetup{colorlinks=true, linkcolor=blue, citecolor=blue, urlcolor=blue}
% long Lean and file names: allow line breaks after dots, slashes and underscores
\makeatletter\g@addto@macro\UrlBreaks{\do\.\do\_}\makeatother
\Urlmuskip=0mu plus 1mu
\emergencystretch=3em
\newcommand{\repo}{https://github.com/mt0-svg/furio-lombardo-quartic}
% robust, so that amsart does not uppercase the address in the author line
\DeclareRobustCommand{\authorlink}{\href{\repo/issues/new/choose}{\nolinkurl{311859884+mt0-svg@users.noreply.github.com}}}


\newtheorem{theorem}{Theorem}[section]
\newtheorem{lemma}[theorem]{Lemma}
\newtheorem{proposition}[theorem]{Proposition}
\newtheorem{corollary}[theorem]{Corollary}
\newtheorem{conjecture}[theorem]{Conjecture}
\theoremstyle{definition}
\newtheorem{definition}[theorem]{Definition}
\newtheorem{remark}[theorem]{Remark}
\numberwithin{equation}{section}

\newcommand{\Q}{\mathbb{Q}}
\newcommand{\Z}{\mathbb{Z}}
\newcommand{\F}{\mathbb{F}}
\newcommand{\Qt}{\mathbb{Q}_2}
\newcommand{\Zt}{\mathbb{Z}_2}
\newcommand{\OK}{\mathcal{O}_K}
\newcommand{\Ov}{\mathcal{O}_v}
\newcommand{\kv}{K_v}
\newcommand{\Gal}{\operatorname{Gal}}
\newcommand{\Cl}{\operatorname{Cl}}
\newcommand{\Res}{\operatorname{Res}}
\newcommand{\pr}{\operatorname{pr}}
\newcommand{\sq}[1]{#1^{\times}/#1^{\times 2}}
% Lean counterpart of a numbered statement: \leanref[note]{declaration}{file}.
\newcommand{\leanref}[3][]{\par\smallskip{\raggedright\noindent\textit{Lean:} \path{#2} in \path{#3}#1.\par}\smallskip}
\newcommand{\nolean}[1]{\par\smallskip{\raggedright\noindent\textit{Lean:} not formalized#1.\par}\smallskip}
\newcommand{\statementfile}{../FurioLombardo/Statement.lean}
\lstset{language={},basicstyle=\linespread{1}\ttfamily\footnotesize,columns=fullflexible,keepspaces=true,
  breaklines=true,breakindent=2em,frame=none,xleftmargin=1em,
  extendedchars=true,inputencoding=utf8,
  literate={ℚ}{{\ensuremath{\mathbb{Q}}}}1 {×}{{\ensuremath{\times}}}1 {•}{{\ensuremath{\bullet}}}1 {→}{{\ensuremath{\to}}}1
    {↔}{{\ensuremath{\leftrightarrow}}}1 {∀}{{\ensuremath{\forall}}}1 {∃}{{\ensuremath{\exists}}}1 {∈}{{\ensuremath{\in}}}1
    {∧}{{\ensuremath{\wedge}}}1 {≠}{{\ensuremath{\neq}}}1 {²}{{\ensuremath{^2}}}1}

\begin{document}

\title[The quartic of Furio and Lombardo]{The plane quartic of Furio and Lombardo has exactly four rational points}
\author{\authorlink}
\thanks{2020 {\em Mathematics Subject Classification.} Primary: 11G30; Secondary: 11D41, 11G10, 11F80, 14H45, 03B35.\\\indent
{\em Key words and phrases.} Rational points; plane quartic; twist of the Klein quartic; Selmer group Chabauty; Prym variety; genus 2 Jacobian; 2-adic logarithm; 7-adic Galois representation; Lean 4.}

\begin{abstract}
Furio and Lombardo conjectured that the plane quartic curve
\[
x^4 + 3x^3y - 3x^2yz - 3x^2z^2 + 6xy^3 - 6xy^2z + 3xyz^2 - 2xz^3 + 4y^4 + 2y^3z - 5yz^3 = 0,
\]
a twist of the Klein quartic that arises in their classification of the $7$-adic images of Galois for elliptic curves over $\mathbb{Q}$, has exactly four rational points. We prove this. A descent over a number field of degree $21$ lifts every rational point to one of two \'etale double covers of the curve. The Prym varieties of these covers are Jacobians of genus $2$ curves, and Stoll's Selmer group Chabauty at one $2$-adic place, applied to them, leaves only the four known points. The proof is formalized in Lean~4 with Mathlib. With the work of Furio and Lombardo, it completes the classification of the $7$-adic images.
\end{abstract}

\maketitle

\definecolor{leangreen}{RGB}{0,122,61}
\definecolor{compblue}{RGB}{0,84,166}
\newcommand{\badge}[2]{\colorbox{#1}{\strut\textcolor{white}{\sffamily\bfseries\small #2}}}
\begin{center}
\small
\renewcommand{\arraystretch}{1.5}
\begin{tabular}{>{\raggedright\arraybackslash}p{0.2\textwidth} >{\raggedright\arraybackslash}p{0.7\textwidth}}
\rowcolor{black!6}
\multicolumn{2}{>{\raggedright\arraybackslash}p{0.93\textwidth}}{\textit{This is AI-generated research: the results, proofs and code were found and written by AI. Credit goes to all the humans whose work it builds on.}}\\
\badge{leangreen}{Proved in Lean 4} & Theorem~\ref{thm:main} is proved in Lean~4 with Mathlib, with no hypothesis and only the standard axioms (\path{FurioLombardo.conjecture_1_6}).\\
\badge{compblue}{By computation} & The certificates that the Lean proof checks were written by PARI/GP scripts and Lean programs, which the repository holds. To guard against a bug in the PARI/GP programs that first computed the local group at $v$ and the covering of $C(\Qt)$, a second program written separately in Sage recomputes the logarithms, the lattice and its saturation, the centres and the boxes, and agrees with them at every centre and every box; it takes the exact data, the points $D_i$, the list of boxes, the image of the Selmer group at $v$ and one matrix of the logarithm from the first programs. Another Sage program redoes the descent of Section~\ref{sec:descent} and finds the same two classes (Section~\ref{sec:comp-second}).\\
\rowcolor{black!6}
\multicolumn{2}{>{\raggedright\arraybackslash}p{0.93\textwidth}}{Lean proofs, code, certificates and the \TeX{} source: \href{\repo}{\texttt{github.com/mt0-svg/furio-lombardo-quartic}}}\\
\end{tabular}
\end{center}

\clearpage
\section{Introduction}\label{sec:intro}

We prove that the plane quartic curve of Conjecture~1.6 of Furio and Lombardo has exactly four rational points. The proof is formalized in Lean~4 with Mathlib, with no hypothesis and only the standard axioms (Section~\ref{sec:comp-lean}).

Furio and Lombardo study the images of the $7$-adic Galois representations attached to elliptic curves over~$\Q$ \cite{FL25}, a case of the classification of Galois images that Mazur proposed as his Program~B \cite[p.~109]{Mazur77}. Following the strategy of Poonen, Schaefer and Stoll for the equation $x^2+y^3=z^7$ \cite{PSS07}, they reduce the open cases of the classification to generalized Fermat equations of signature $(2,3,7)$, and these to the rational points of finitely many twists of the Klein quartic. They determine the rational points of all of these twists but one, the curve they call $X_{E_3}$, for which they state the following conjecture. We quote it from the arXiv version of 2026-03-05 (v3), where it has the same number as in the journal version.

\begin{conjecture}[{Furio and Lombardo \cite[Conjecture~1.6]{FL25}}]\label{conj:FL}
The set of rational points of the curve
\[
C:\ x^4 + 3x^3y - 3x^2yz - 3x^2z^2 + 6xy^3 - 6xy^2z + 3xyz^2 - 2xz^3 + 4y^4 + 2y^3z - 5yz^3 = 0
\]
is $C(\Q)=\{[0:0:1],[1:1:1],[2:0:1],[-1:0:1]\}$.
\end{conjecture}

\begin{theorem}\label{thm:main}
Let $F(x,y,z)$ be the quartic form of Conjecture~\ref{conj:FL}. For all rational numbers $x,y,z$, not all zero, $F(x,y,z)=0$ holds if and only if $(x,y,z)$ is a nonzero rational multiple of one of the four vectors $(0,0,1)$, $(1,1,1)$, $(2,0,1)$, $(-1,0,1)$. In particular Conjecture~\ref{conj:FL} is true.
\end{theorem}
\leanref[; the statement \path{FurioLombardo.Conjecture} is in \path{FurioLombardo/Statement.lean} (Section~\ref{sec:comp-lean})]{FurioLombardo.conjecture_1_6}{FurioLombardo/Main.lean}

The four points lie on $C$, and a search over the primitive integer vectors with coordinates of absolute value at most $100$ finds no other point. The following lines of the file \path{paper/small_checks.gp} of the repository do both in a few seconds with PARI/GP \cite{PARI},
\lstinputlisting[linerange={2-7}]{small_checks.gp}
and print
\lstinputlisting[linerange={1-2}]{small_checks.out}
the second line being the list of the points found, one vector for each.

Theorem~\ref{thm:main} proves Conjecture~\ref{conj:FL}. Furio and Lombardo prove three results under this conjecture: the classification of the $7$-adic images of Galois for elliptic curves over $\Q$ without complex multiplication \cite[Theorem~1.7]{FL25}, the list of the solutions with property $(\star)$ of the equation $a^2+196b^3=27c^7$ \cite[Theorem~3.5(2)]{FL25}, and the statement that an elliptic curve over $\Q$ whose image modulo $49$ lies in one of the groups $G_{ns}^\#(49)$ and $G_{sp}^\#(49)$ has $j$-invariant~$0$ \cite[Corollary~3.7(2)]{FL25}. With Theorem~\ref{thm:main} these three results hold unconditionally; this deduction is theirs, and it is not formalized here. In \cite[Theorem~3.5(2)]{FL25} the solutions are printed as $(\pm13,196,-1)$; they are $(\pm13,-1,-1)$, as found by the search of \cite[Remark~3.6]{FL25}, since $13^2+196\cdot(-1)^3=-27=27\cdot(-1)^7$, as Nguyen Xuan Tho notes \cite{NguyenXuanTho26}.

\subsection*{What was known}
Furio and Lombardo show that the Jacobian $J$ of $C$ has rank $3$ over $\Q$ \cite[Remark~6.3]{FL25}, which is the genus of $C$, so the method of Chabauty and Coleman does not apply \cite[Section~6.5]{FL25}. They note that the bitangents of $C$ are defined over fields of degree $28$, which makes a descent by the corresponding \'etale double covers impractical, and they suggest quadratic Chabauty over number fields, as used by Balakrishnan, Betts, Hast, Jha and M\"uller for the non-split Cartan modular curve of level $27$ \cite{BBHJM25}. Two independent proofs of Conjecture~\ref{conj:FL} were claimed shortly before this paper, both through a map of degree two from $C$ to an elliptic curve over a field of degree~$21$. In a one-page abstract of 2026-09-20, Karabiyik announces a proof that combines this quotient with $37$-adic height equations \cite{Karabiyik26}; the paper is in preparation. On 2026-09-28, Nguyen Xuan Tho posted a preprint \cite{NguyenXuanTho26} in which the elliptic curve has complex multiplication and the proof combines a Mordell-Weil sieve with quadratic Chabauty at the prime $5581$; it takes the rank~$3$ of $J(\Q)$ from \cite[Remark~6.3]{FL25}. Neither comes with a formal proof. The proof below uses another method and does not rely on them.

\subsection*{The method}
Let $K$ be the field of definition of a $2$-torsion point of $J$ whose Galois orbit has $21$ elements, a field of degree~$21$. Bruin \cite{Bruin08} attaches to such a point quadratic forms with $Q_1Q_3-Q_2^2=c_BF$ and, for each $\delta\in K^\times$, an \'etale double cover $D_\delta\to C$ over $K$ whose Prym variety is the Jacobian of a genus $2$ curve. A descent over $K$, which uses the class number of $K$ and local conditions at finitely many primes, shows that every rational point of $C$ lifts to $D_\delta(K)$ for one of two classes $\delta$ (Section~\ref{sec:descent}). The proof uses no value of the rank of $J(\Q)$: this descent replaces the Mordell-Weil sieve (Remark~\ref{rem:alt-descent}). For each of the two covers we then apply Stoll's Selmer group Chabauty \cite[Theorem~2.1]{Stoll17} at the place $v$ of $K$ above $2$ with ramification index~$3$, to the group $A_k(K)$ of $K$-points of the genus $2$ Jacobian (Sections~\ref{sec:prym} to~\ref{sec:proof}). This works for two reasons. A $2$-descent over $K$ shows that $A_k(K)/2A_k(K)$ injects into $A_k(\kv)/2A_k(\kv)$ (Corollary~\ref{cor:selmer-consequences}), so the argument needs neither generators of $A_k(K)$ nor its rank. And $C(\Qt)$ is covered by finitely many boxes, each of which either contains no point that lifts to the cover over $\kv$ or carries a $2$-adic logarithm with a constant leading class, which checks the remaining condition of Stoll's criterion at every $2$-adic point (Section~\ref{sec:cover}).

\subsection*{The formalization}
Each numbered statement below is followed by the name of the Lean declaration that proves it and of its file, or by the words ``not formalized''. The finite computations, from the class number certificate to the boxes of the covering, are checked by the Lean kernel, and no compiled code is trusted. Section~\ref{sec:comp-lean} prints the Lean statement of Theorem~\ref{thm:main} and describes these checks.

Section~\ref{sec:descent} carries out the descent over $K$. Section~\ref{sec:prym} introduces the genus $2$ curves, a model of the points of their Jacobians and the Abel-Prym map. Section~\ref{sec:stoll} proves Stoll's criterion at one place in the form used here, Section~\ref{sec:selmer} the Selmer bound, and Section~\ref{sec:local} describes the local group at $v$ through its logarithm. Section~\ref{sec:cover} covers $C(\Qt)$, and Section~\ref{sec:proof} assembles the proof. Section~\ref{sec:comp} describes the computations and their checks. The notation of Section~\ref{sec:descent} is fixed throughout, and $k\in\{0,1\}$ indexes the two covers.

\section{The descent over a field of degree 21}\label{sec:descent}

This section attaches to every rational point of $C$ a square class in a number field $K$ of degree~$21$ and proves that only two classes occur (Theorem~\ref{thm:selmer-set}). The argument uses exact arithmetic in $K$, the class number of $K$, and local conditions at finitely many primes; it does not use the Jacobian of~$C$.

\subsection{The field}
Let $g(x)\in\Z[x]$ be the polynomial
\[
\begin{aligned}
g(x)={}&x^{21}-7x^{20}+14x^{19}-84x^{16}+98x^{15}+2x^{14}+175x^{13}-609x^{12}+980x^{11}-770x^{10}-280x^{9}\\
&+1008x^{8}-1072x^{7}+560x^{6}+28x^{5}-336x^{4}+252x^{3}-112x^{2}+28x-4,
\end{aligned}
\]
let $K=\Q(\beta)$ with $g(\beta)=0$, and let $\OK$ be its ring of integers. The polynomial $g$ is irreducible, $K$ has three real and nine complex places, and its discriminant is $-2^{22}\cdot7^{27}$. The prime $2$ has three places in~$K$, all of residue degree one, with ramification indices $3$, $6$ and $12$; we call them $v$, $w_6$ and $w_{12}$. The prime $7$ has one place $w_7$, with ramification index $7$ and residue degree $3$. The field $K$ is the field of definition of a $2$-torsion point of the Jacobian of $C$ whose Galois orbit has $21$ elements; this origin explains the arithmetic below but is not used in the proofs.

\begin{proposition}\label{prop:classnumber}
The ring $\OK$ is a principal ideal domain.
\end{proposition}
\leanref[; with no hypothesis, and on the field of Section~\ref{sec:descent} through \path{FurioLombardo.M2.Proved.clK21TwoTorsionTrivial}]{FurioLombardo.M2.Proved.isPrincipalIdealRing_K21}{FurioLombardo/M2/ZimmertProved.lean}

\begin{proof}
We use Zimmert's bound \cite[Satz~2]{Zimmert81} in the case $\gamma=\tfrac12$. For a number field with $r_1$ real and $r_2$ complex places and $0<\alpha<\tfrac12$, put $\psi=\Gamma'/\Gamma$ and
\[
\begin{aligned}
S_2(r_1,r_2;\alpha)={}&r_1\Bigl(-\psi\bigl(\tfrac34\bigr)+\log\Gamma\bigl(\tfrac32\bigr)+\tfrac12\log\pi\Bigr)
+r_2\Bigl(-2\psi\bigl(\tfrac32\bigr)+2\log2+\log\pi\Bigr)\\
&-\frac{2}{\frac12-\alpha}-\log\Bigl(\frac{(1+\alpha^{-1})}{9\,\bigl(1+(1-\alpha)^{-1}\bigr)}\Bigr).
\end{aligned}
\]
Zimmert's theorem says that every ideal class contains a nonzero integral ideal $I$ with $S_2(r_1,r_2;\alpha)\le\log\bigl(\sqrt{|d_K|}/N(I)\bigr)$, where $d_K$ is the discriminant. For $K$, with $(r_1,r_2)=(3,9)$ and $\alpha=\tfrac1{12}$, the value $S_2(3,9;\tfrac1{12})=21\gamma_E-\tfrac{3\pi}{2}+60\log2+12\log\pi-\tfrac{204}{5}+\log\tfrac{207}{143}\approx22.3046$, with $\gamma_E$ Euler's constant, exceeds $\log\sqrt{2^{22}7^{27}}-\log120000\approx22.1992$. Since $|d_K|\le2^{22}7^{27}$, every ideal class of $K$ contains an integral ideal of norm at most $120000$, and it suffices that every prime ideal of norm at most $120000$ is principal. Let $\mathfrak p$ be such a prime, above the rational prime~$p$. If $p\notin\{2,7,45613\}$, then, since $p\le120000$, $p$ does not divide the resultant of $g$ and $g'$, which lies in the conductor of $\Z[\beta]$ in $\OK$, so $\mathfrak p=(p,\phi(\beta))$ for a monic irreducible factor $\phi$ of $g$ modulo~$p$ (Dedekind and Kummer). The certificate lists the factors of $g$ modulo $p$ of each degree $e$ with $p^e\le120000$, which it checks through the identity of $\prod\phi$ with $\gcd(g,x^{p^e}-x)$, and for each of them an element $\alpha\in\OK$ and an element $c\in\OK$ with $\alpha c=p$, $c\notin\mathfrak p$ and $\phi(\beta)c\in p\,\OK$; then $\mathfrak p=(\alpha)$. There is one such certificate for each prime $p\le120000$ other than $2$, $7$ and $45613$. The primes above $2$ are generated by factors of $\beta$ and $\beta-1$, the prime above $7$ by an element $\alpha_7$ with $\alpha_7^7\varepsilon=7$ for a unit $\varepsilon$, and every prime of degree one above $45613$ contains an explicit element of norm $\pm45613$, which generates it.
\end{proof}

The Minkowski bound of $K$ is about $4.04\cdot10^{7}$; the smaller bound $120000$ is what makes the certificate short. The class number was first computed with the Minkowski bound in PARI/GP (Section~\ref{sec:comp}).

\subsection{The descent class}

\begin{lemma}\label{lem:bruin-identity}
There are ternary quadratic forms $Q_1,Q_2,Q_3$ with coefficients in $\OK$ and an element $c_B\in\OK$ such that
\begin{equation}\label{eq:bruin}
Q_1Q_3-Q_2^2=c_B\,F .
\end{equation}
\end{lemma}
\leanref{FurioLombardo.M1.bruin_O}{FurioLombardo/M1/BruinCheck.lean}

The coefficients of the forms and of $c_B$ are in the repository; the coordinates of $c_B$ in the integral basis of $\OK$ are not all zero, so $c_B\neq0$. The forms come from the $2$-torsion point of the Jacobian of $C$ over $K$ mentioned above, by the construction of Bruin \cite[Lemma~3.3 and Section~5, Case~4]{Bruin08}; the proofs below only use the identity~\eqref{eq:bruin}. For a rational point $P=[a:b:c]$ of $C$ put $q(P)=Q_1(a,b,c)$ if this is nonzero and $q(P)=Q_3(a,b,c)$ otherwise. By Lemma~\ref{lem:support} below, $q(P)\neq0$. Its class in $\sq{K}$ does not depend on the choice of homogeneous coordinates, since the forms are quadratic, and we call it the descent class of~$P$.

Let $S$ be the set of the six places $v$, $w_6$, $w_{12}$, $w_7$, $\mathfrak q_3$ and $\mathfrak q_{439}$ of $K$, where $\mathfrak q_3$ and $\mathfrak q_{439}$ are places of residue degree one above $3$ and $439$.

\begin{lemma}\label{lem:support}
Let $(a,b,c)$ be a primitive integer vector with $F(a,b,c)=0$, let $P=[a:b:c]$, and let $\mathfrak p$ be a prime of $\OK$ outside~$S$. Then $Q_1(a,b,c)$, $Q_2(a,b,c)$, $Q_3(a,b,c)$ do not all lie in $\mathfrak p$. Consequently $q(P)\neq0$, and the valuation $v_{\mathfrak p}(q(P))$ is even.
\end{lemma}
\leanref[; $q(P)\neq0$ is \path{FurioLombardo.M1.exists_qO_ne_zero} in \path{FurioLombardo/M1/Sieve.lean}, and the parity is \path{FurioLombardo.M1.mem_selmer_of_qO} in \path{FurioLombardo/M1/Descent.lean}]{FurioLombardo.M1.not_all_mem}{FurioLombardo/M1/Descent.lean}

\begin{proof}
Let $J=\det(\partial Q_i/\partial x_j)$, a ternary cubic form, and let $M_6$ be the $6\times6$ matrix whose rows are the coefficient vectors of $Q_1,Q_2,Q_3,\partial J/\partial x,\partial J/\partial y,\partial J/\partial z$ in the monomial basis $x^2,xy,xz,y^2,yz,z^2$. Let $R$ be a point of $k^3$, for a field $k$, at which the reductions of $Q_1,Q_2,Q_3$ vanish. Suppose $R\neq0$ and let $a$ be an index with $R_a\neq0$. By Euler's identity, $R_aJ(R)$ is twice the determinant of the Jacobian matrix at $R$ with its column $a$ replaced by the values of the forms, so $J(R)=0$. A second identity writes $R_a(\partial J/\partial x_l)(R)$ as twice a sum of determinants that have the values of the forms as column $a$, plus the determinant of the Jacobian matrix with its column $a$ replaced by its column $l$, which is $J(R)$ if $l=a$ and $0$ otherwise; so $(\partial J/\partial x_l)(R)=0$ for each~$l$. Both identities are polynomial identities with integer coefficients, so this holds in every characteristic. The six forms then vanish at $R$, so $M_6$ kills the nonzero vector of quadratic monomials of~$R$. We checked the identity $N\cdot M_6=c_S\cdot I_6$ over $\OK$ for an explicit matrix $N$, with $c_S=2^{13}7^{13}\pi_3^{12}\pi_{439}^{12}$, where $\pi_3$ and $\pi_{439}$ generate $\mathfrak q_3$ and $\mathfrak q_{439}$. Since $c_S\notin\mathfrak p$, the reductions of $Q_1,Q_2,Q_3$ modulo $\mathfrak p$ have no common zero other than $0$, and the reduction of $(a,b,c)$ is nonzero; this is the first claim. On $C$ the identity~\eqref{eq:bruin} reads $Q_1Q_3=Q_2^2$. Applied to a prime above $5$, which is outside $S$, the first claim shows that $Q_1(a,b,c)$ and $Q_3(a,b,c)$ are not both zero, since otherwise $Q_2(a,b,c)^2=0$ as well; so $q(P)\neq0$. If $Q_1(a,b,c)$ and $Q_3(a,b,c)$ are both nonzero, one of them is a unit at $\mathfrak p$, because $Q_2(a,b,c)^2=Q_1Q_3$ lies in $\mathfrak p$ when both do, and $v_{\mathfrak p}(Q_1)+v_{\mathfrak p}(Q_3)=2v_{\mathfrak p}(Q_2)$ then shows that both valuations are even. If one of them is zero, then $Q_2(a,b,c)=0$, and the other one is a unit at~$\mathfrak p$.
\end{proof}

Let $K(S,2)$ be the subgroup of $\sq{K}$ of classes with even valuation at every prime outside~$S$.

\begin{lemma}\label{lem:sunits-K}
The group $K(S,2)$ has dimension $18$ over $\F_2$. It has the basis formed by the classes of $-1$, of eleven explicit units and of generators of the six prime ideals of~$S$.
\end{lemma}
\leanref[; the independence is \path{FurioLombardo.M1.gProd_eq_sq} in \path{FurioLombardo/M1/Sieve.lean}]{FurioLombardo.M1.exists_isSquare_mul_prod}{FurioLombardo/M1/SelmerBound.lean}

\begin{proof}
Since the class group of $K$ has odd order (Proposition~\ref{prop:classnumber}), an element with even valuations outside $S$ is an $S$-unit times a square, so $K(S,2)$ is the image of the $S$-units. The unit group has rank $r_1+r_2-1=11$ and torsion $\{\pm1\}$, and $\#S=6$, so the image has dimension at most $1+11+6=18$ and is spanned by the listed elements. They are independent: we computed $32$ quadratic residue characters of the $18$ elements, at primes of $K$ above $5$, $11$, $13$, $17$, $19$ and $23$, and the resulting $32\times18$ matrix over $\F_2$ has a left inverse.
\end{proof}

\begin{lemma}\label{lem:sieve}
There is an explicit set of eight classes of $K(S,2)$ that contains the descent class of every rational point of~$C$. It is cut out by the reductions of $C$ modulo $5$, $11$, $13$, $17$, $19$ and $23$.
\end{lemma}
\leanref{FurioLombardo.M1.sieve_good}{FurioLombardo/M1/Sieve.lean}

\begin{proof}
Let $\mathfrak p$ be one of the $32$ primes of the proof of Lemma~\ref{lem:sunits-K}, above a prime $p$, and let $\chi_{\mathfrak p}$ be the quadratic character of $(\OK/\mathfrak p)^\times$. Let $(a,b,c)$ be a primitive integer vector of $P$. By Lemmas~\ref{lem:support} and~\ref{lem:sunits-K} there is a product $G$ of basis elements such that $Q_i(a,b,c)\,G$ is a square in $\OK$ for each $i\in\{1,3\}$ with $Q_i(a,b,c)\neq0$; the basis elements are units at~$\mathfrak p$. Let $\bar P$ be the reduction of $(a,b,c)$ modulo~$p$, a point of $C(\F_p)$. By Lemma~\ref{lem:support} and the identity $Q_1Q_3=Q_2^2$ on $C$, one of $Q_1(\bar P)$, $Q_3(\bar P)$ is nonzero in $\OK/\mathfrak p$, and $\chi_{\mathfrak p}(G)=\chi_{\mathfrak p}(Q_i(\bar P))$ for that~$i$. So the exponent vector of $G$ over $\F_2$ satisfies, for each of the $32$ primes $\mathfrak p$, one of the affine conditions given by the points of $C(\F_p)$. We enumerated $C(\F_p)$ for the six primes. The points give $128$ possible vectors of values of the $32$ characters, each of them determines the exponent vector through the left inverse of Lemma~\ref{lem:sunits-K}, and whenever that exponent vector reproduces the character vector it is one of eight explicit vectors.
\end{proof}

\begin{lemma}\label{lem:local-kills}
Four of the eight classes of Lemma~\ref{lem:sieve} are excluded by the points of $C$ modulo $27$, read at the place $\mathfrak q_3$, and two more by the points of $C$ modulo $8$, read at the place~$w_6$. The two remaining classes are the classes $\delta_0$ and $\delta_1$ of $q(P_0)=q(P_2)$ and of $q(P_1)=q(P_3)$, where $P_0=[0:0:1]$, $P_1=[1:1:1]$, $P_2=[2:0:1]$ and $P_3=[-1:0:1]$.
\end{lemma}
\leanref[ and \path{FurioLombardo.M1.kill2} in \path{FurioLombardo/M1/Local2.lean}; the two remaining classes are \path{FurioLombardo.M1.ck_known0} and \path{FurioLombardo.M1.ck_known1} with \path{FurioLombardo.M1.survOf} in \path{FurioLombardo/M1/SelmerSet.lean}]{FurioLombardo.M1.kill3}{FurioLombardo/M1/Local3.lean}

\begin{proof}
Let $(a,b,c)$ be a primitive integer vector of a rational point, and let $G$ be as in the proof of Lemma~\ref{lem:sieve}, a product of basis elements such that $Q_i(a,b,c)G$ is a square in $\OK$ for each $i\in\{1,3\}$ with $Q_i(a,b,c)\neq0$. The polynomial $g$ has the root $17$ modulo $27$, which gives a ring homomorphism $\Z[\beta]\to\Z/27\Z$ sending $\beta$ to~$17$. Since $3$ does not divide the discriminant of $g$, it does not divide the index of $\Z[\beta]$ in $\OK$, and the homomorphism extends to~$\OK$. Some coordinate of $(a,b,c)$ is prime to $3$, so modulo $27$ the vector is a unit times a point of one of the planes $x=1$, $y=1$, $z=1$. For each of the $3\cdot27^2$ points of these planes, either $F$ does not vanish modulo $27$, or for each of four of the eight classes one of $Q_1G$, $Q_3G$ is not a square modulo $27$; this is a finite check. Since squares of $\OK$ map to squares modulo $27$, and $0$ is a square, it excludes these four classes.

The generator $\lambda$ of $w_6$ has norm $\pm2$, and $2=\lambda^6U$ with $U$ a unit at $w_6$, so $8=\lambda^{18}U^3$. The vector $(a,b,c)$ is $l\,P+8b'$ with $l$ odd, $b'$ integral and $P$ one of the $32$ points modulo $8$ of the charts $(x,y,1)$, $(1,y,2z)$, $(2x,1,2z)$ at which $F\equiv0\pmod 8$. For each of these points and each of two further classes, with $G\equiv G'\pmod 8$, the certificate names $i\in\{1,3\}$ and writes $w=Q_i(P)G'$ either as $\lambda^{m}(1+\lambda t)$ with $m$ odd and $m<18$, or as $\lambda^{2j}\bigl((1+\lambda s)^2+\lambda^{d}(1+\lambda z)\bigr)$ with $d$ odd, $d<12$ and $2j+d<18$, for explicit $t,s,z\in\OK$. Then $Q_i(a,b,c)G=l^2w+8r$ with $r\in\OK$, and in both cases this element is not a square: in the first case its valuation at $w_6$ is odd, in the second case it is $\lambda^{2j}$ times a unit of the form $(1+\lambda s)^2+\lambda^{d}u$ with $u$ a unit and $d$ odd and less than $12=2v_{w_6}(2)$, which is not a square. This contradicts the choice of $G$, also when $Q_i(a,b,c)=0$, since $0$ is a square, and excludes the two classes.

For the two remaining representatives $G_0$ and $G_1$, the products $\delta_0G_0$ and $\delta_1G_1$ are explicit squares in~$K$. The values $q(P_i)$ at the four known points are computed directly.
\end{proof}

\begin{theorem}\label{thm:selmer-set}
For every rational point $P$ of $C$, the descent class of $P$ is $\delta_0$ or $\delta_1$: there are $k\in\{0,1\}$ and $t\in K$ with $q(P)=\delta_kt^2$.
\end{theorem}
\leanref[; with Proposition~\ref{prop:classnumber} as input through \path{FurioLombardo.M2.Proved.clK21TwoTorsionTrivial}; the main theorem uses the form at a primitive integer vector, \path{FurioLombardo.M1.selmer_int} in the same file]{FurioLombardo.M1.selmerSet}{FurioLombardo/M1/SelmerSet.lean}

\begin{proof}
Lemma~\ref{lem:support} puts the descent class in $K(S,2)$, and Lemmas~\ref{lem:sieve} and~\ref{lem:local-kills} leave $\delta_0$ and $\delta_1$.
\end{proof}

We fix representatives $\delta_0,\delta_1\in K^\times$ of the two classes. They are units at the three places above~$2$, and $\delta_0/\delta_1$ is not a square. For $\delta\in K^\times$ let $D_\delta\subset\mathbb P^4$ be the curve with coordinates $(x:y:z:r:s)$ defined by
\begin{equation}\label{eq:Ddelta}
Q_1(x,y,z)=\delta r^2,\qquad Q_2(x,y,z)=\delta rs,\qquad Q_3(x,y,z)=\delta s^2,
\end{equation}
with the involution $\iota(x:y:z:r:s)=(x:y:z:-r:-s)$ and the projection $\pi(x:y:z:r:s)=(x:y:z)$. On the image of $\pi$, the identity~\eqref{eq:bruin} gives $c_BF=Q_1Q_3-Q_2^2=0$, so $\pi$ maps $D_\delta$ to~$C$.

\begin{corollary}\label{cor:lift}
Every rational point $P$ of $C$ is the image under $\pi$ of a point of $D_{\delta_0}(K)$ or of $D_{\delta_1}(K)$. The points $P_0,P_2$ lift to points $x_0,x_2$ of $D_{\delta_0}(K)$ and the points $P_1,P_3$ lift to points $x_1,x_3$ of $D_{\delta_1}(K)$.
\end{corollary}
\leanref[; with Proposition~\ref{prop:classnumber} as input through \path{FurioLombardo.M2.Proved.clK21TwoTorsionTrivial}; the lifts $x_0,\dots,x_3$ are \path{FurioLombardo.Discharge.M3a.Bruin.x0} to \path{FurioLombardo.Discharge.M3a.Bruin.x3}, with their images \path{x0_over} to \path{x3_over}, in \path{FurioLombardo/Discharge/M3a/Concrete.lean}]{FurioLombardo.M1.descent_Mmat}{FurioLombardo/M1/DescentM3a.lean}

\begin{proof}
This is the lifting lemma of Bruin \cite[Lemma~8.1]{Bruin08}, with the two classes of Theorem~\ref{thm:selmer-set}. If $Q_1(P)=\delta_kt^2\neq0$, the point $(P:t:Q_2(P)/(\delta_kt))$ lies on $D_{\delta_k}$, since $\delta_k(Q_2/(\delta_kt))^2=Q_2^2/Q_1=Q_3$ on $C$. If $Q_1(P)=0$, then $Q_2(P)=0$ and $Q_3(P)=\delta_kt^2$, and $(P:0:t)$ lies on $D_{\delta_k}$.
\end{proof}

\begin{remark}\label{rem:alt-descent}
Theorem~\ref{thm:selmer-set} can also be derived from the Mordell-Weil group of the Jacobian $J$ of $C$ over~$\Q$. Furio and Lombardo show that $J(\Q)$ has rank~$3$ \cite[Remark~6.3]{FL25}; a $2$-descent in PARI/GP that uses neither Magma nor the generalized Riemann hypothesis shows that the images of two explicit divisor classes span the fake $2$-Selmer group, of dimension~$2$, and that the kernel of the descent map on $J(\Q)/2J(\Q)$ is spanned by a third explicit class, so that these three classes form a basis of $J(\Q)/2J(\Q)$. The torsion of $J(\Q)$ injects into $J(\F_3)$ and into $J(\F_{11})$, whose orders $40$ and $3277$ are coprime, so $J(\Q)$ has no torsion, and the $2$-descent gives the rank~$3$ again. A Mordell-Weil sieve shows that every rational point $P$ satisfies $[P-P_i]\in2J(\Q)$ for some~$i$. The divisor of the function $Q_1/l^2$, for a linear form $l$, is twice a divisor, so by Weil reciprocity this function induces a homomorphism $J(K)\to\sq{K}$, which sends $[P-P_i]$ to the class of $q(P)/q(P_i)$; hence $q(P)$ has the class of $q(P_i)$. The proof of this paper and its Lean formalization use no value of the rank of $J(\Q)$: the descent over $K$ of this section replaces the Mordell-Weil sieve.
\end{remark}
\nolean{; the $2$-descent and the sieve are checked by computation, Section~\ref{sec:comp-other}}

\section{The genus 2 curves and the Abel-Prym map}\label{sec:prym}

This section introduces, for each of the two twists, a genus~$2$ curve over $K$, a model of the group of points of its Jacobian that works over every field containing $K$, and a map from $D_\delta(K)$ to that group.

\subsection{The curves}
Write $Q_i(p)=p^{\mathsf T}M_ip$ with symmetric $3\times3$ matrices $M_i$ over $K$, and for $\delta\in K^\times$ put
\[
f_\delta(t)=-\delta\det\bigl(M_1+2tM_2+t^2M_3\bigr)\in K[t].
\]
Bruin \cite[Section~8 and Theorem~5.1]{Bruin08} shows that the Prym variety of the double cover $D_\delta\to C$ is the Jacobian of the genus $2$ curve $Y^2=f_\delta(t)$. We do not use this identification: Theorem~\ref{thm:stoll-one-place} below applies to any map from the points of $D_\delta$ to an abelian group. We work with the reversed polynomials
\[
f_k(X)=X^6f_{\delta_k}(1/X)\qquad(k=0,1),
\]
because the leading coefficient of $f_{\delta_0}$ is a square in $\kv$, while the leading coefficient $f_{\delta_k}(0)$ of $f_k$ is not a square in $\kv$, for both~$k$.

Let $L$ be a field of characteristic different from~$2$. We call $f\in L[X]$ a good sextic if it is squarefree of degree~$6$ and its leading coefficient is not a square in~$L$.

\begin{lemma}\label{lem:fk}
For $k=0,1$, the polynomial $f_k$ has degree $6$ and is a good sextic over $K$ and over~$\kv$. It factors over $K$ as $f_k=c_k\,q_k\,h_k$ with $c_k\in K^\times$ and $q_k$, $h_k$ monic of degrees $2$ and $4$. The discriminant $d_k$ of $q_k$ is not a square in $\kv$, and $h_k$ is irreducible over~$\kv$.
\end{lemma}
\leanref[; over $\kv$ it is \path{goodSextic_fRev_Kv} in \path{FurioLombardo/Discharge/M3a/LocalLead.lean}, the local facts are \path{d_not_isSquare_Kv} and \path{irr_Kv} in \path{FurioLombardo/Discharge/M3a/LocalFacts.lean}, the factorization is \path{fRev_eq_mul}]{FurioLombardo.Discharge.M3a.Bruin.goodSextic_fRev}{FurioLombardo/Discharge/M3a/Concrete.lean}

\begin{proof}
The coefficients of $f_k$ and of its factors are explicit, and the identities $f_k=c_kq_kh_k$ and $f_k=X^6f_{\delta_k}(1/X)$, with $f_{\delta_k}$ given by the determinant, are checked by exact computation. A B\'ezout identity $af_k+bf_k'=1$ with explicit $a,b\in K[X]$ shows that $f_k$ is squarefree. The leading coefficient $c_k$ of $f_k$ is a quadratic nonresidue modulo a prime of degree one above $13$ or $23$, so it is not a square in $K$, and the residue of $4c_k$ modulo $8$ in the ring of integers of $\kv$ is not the residue of a square, so it is not a square in~$\kv$. Over $N=\kv(\omega)$ with $\omega^2=d_k$, the quartic $h_k$ is $\Psi\bar\Psi$ with $\Psi=X^2+(a_1+b_1\omega)X+(a_0+b_0\omega)$ and $\bar\Psi$ its conjugate, for explicit $a_i,b_i\in K$ with $b_1\neq0$; it is irreducible over $\kv$ because the discriminant of $\Psi$ is not a square in~$N$. The two nonsquare statements, for $d_k$ in $\kv$ and for that discriminant in $N$, are proved from explicit $2$-adic approximations to a precision at which the square class is decided.
\end{proof}

\subsection{The group}\label{sec:group}
Let $f\in L[X]$ be a good sextic and $R_f=L[X,Y]/(Y^2-f)$. The ring $R_f$ is a Dedekind domain. The norm of a nonzero ideal $I$ of $R_f$ to $L[X]$ is a principal ideal $(n_I)$, and the parity of $\deg n_I$ is additive and vanishes on principal ideals, because the norm $p^2-q^2f$ of an element $p+qY$ has even degree when the leading coefficient of $f$ is not a square. So the parity defines a homomorphism from the class group $\Cl(R_f)$ to $\Z/2\Z$, and we put
\[
A_f(L)=\ker\bigl(\Cl(R_f)\to\Z/2\Z\bigr).
\]
For the curve $\mathcal F:Y^2=f(X)$, the group $A_f(L)$ is the group of $L$-rational points of the Jacobian of~$\mathcal F$: the two points at infinity of $\mathcal F$ are conjugate over $L$, so they form one place $\infty$ of degree~$2$, $\Cl(R_f)$ is the Picard group of $\mathcal F$ modulo the class of $\infty$, and the kernel of the parity is the image of the divisor classes of degree~$0$. This identification is not used; all statements below are about $A_f(L)$. We write $A_f(L)$ additively. Every class of $A_f(L)$ is trivial or is the class of a Mumford ideal $\langle u,Y-w\rangle$ with $u$ monic of degree~$2$, $\deg w\le1$ and $u\mid w^2-f$, and the class of $\langle u,Y+w\rangle$ is its opposite. A field embedding $\sigma:L\to L'$ induces a homomorphism $A_f(L)\to A_{\sigma f}(L')$ on ideal classes.
\leanref[; the reduced form is \path{jac_eq_one_or_reduced} in \path{FurioLombardo/Discharge/M3a/Reduction.lean}, the base change \path{jacMap} in \path{FurioLombardo/M3a/BaseChange.lean}]{FurioLombardo.M3a.Genus2.Jac}{FurioLombardo/M3a/Pic.lean}

For a monic quadratic factor $q$ of $f$ let $T_q$ be the class of $\langle q,Y\rangle$.

\begin{lemma}\label{lem:two-torsion-T}
The class $T_q$ lies in $A_f(L)$ and satisfies $2T_q=0$.
\end{lemma}
\leanref{FurioLombardo.M3a.Genus2.Tpt_sq}{FurioLombardo/M3a/TwoTorsion.lean}

Let $\mathcal A_f=L[X]/(f)$, with $\Theta$ the class of $X$, and $H_f=\mathcal A_f^\times/L^\times\mathcal A_f^{\times2}$. The $x-T$ map \cite{Schaefer95,Stoll01} is the homomorphism $\mu:\Cl(R_f)\to H_f$ that sends the class of a nonzero ideal $I$ prime to $(Y)$ to the class of $n_I(\Theta)$; every class contains such an ideal, and the value does not depend on the choice. On a Mumford ideal with $u$ prime to $f$ it gives the class of $u(\Theta)$. It is compatible with field embeddings, and it kills $2A_f(L)$ because $H_f$ has exponent~$2$.
\leanref[, on $A_f(L)$ \path{muJ}, with base change \path{muJ_jacMap}]{FurioLombardo.M3a.Genus2.mu}{FurioLombardo/M3a/XminusT.lean}

\begin{proposition}\label{prop:xT-kernel}
For $k=0,1$ the kernel of $\mu:A_{f_k}(K)\to H_{f_k}$ is $2A_{f_k}(K)$.
\end{proposition}
\leanref[, with the hypothesis discharged by \path{poonenSchaefer_fRev} in \path{FurioLombardo/Discharge/M3a/K1Concrete.lean}; the general statement is \path{FurioLombardo.Discharge.M3a.xtKernel_of_deltaTrivial} in \path{FurioLombardo/Discharge/M3a/K1Kernel.lean}]{FurioLombardo.Discharge.M3a.Bruin.xtKernel_fRev}{FurioLombardo/Discharge/M3a/ConcreteInputs.lean}

\begin{proof}
This is the elementary descent of Cassels and Flynn \cite[Chapter~6]{CF96}, carried out in the model $A_f(L)$. We first prove the following statement. Let $f$ be a good sextic over a field $L$ such that
\begin{enumerate}
\item[(H1)] some $d\in L^\times$ and $\varepsilon\in\mathcal A_f$ satisfy $\varepsilon^2=d$ and $N_{\mathcal A_f/L}(\varepsilon)=-d^3$;
\item[(H2)] $f$ has no root in $L$;
\item[(H3)] $\mu(T_{q'})\neq0$ for every monic quadratic factor $q'$ of~$f$.
\end{enumerate}
Then every $Q\in A_f(L)$ with $\mu(Q)=0$ lies in $2A_f(L)$. We may assume that $Q$ is the class of a Mumford ideal $\langle u,Y-w\rangle$. If $u$ and $f$ had a common factor, then by (H2) $u$ would divide $f$ and $w$ would vanish, so $Q=T_u$, which (H3) excludes; so $u$ is prime to $f$ and $u(\Theta)=\eta\rho^2$ with $\eta\in L^\times$ and $\rho\in\mathcal A_f^\times$. Taking norms and using resultants gives $f_6\,\eta^3N(\rho)=\pm N_{L[X]/(u)}(w)$, with $f_6$ the leading coefficient of~$f$. Replacing $(\eta,\rho)$ by $(\eta/d,\varepsilon\rho)$ changes the sign, by (H1). With the sign $-$, the explicit formulas of \cite[Chapter~6, Sections~1, 2, 5 and~8]{CF96} give a Mumford ideal whose class $R$ satisfies $2R=Q$.

For $f_k$ over $K$: the discriminant $d_k$ of $q_k$ is not a square in $K$ and is a square in $\mathcal A_{f_k}$, by an explicit identity $\beta_k^2-d_k=f_k\gamma_k$ in $K[X]$; with $\varepsilon$ the class of $\beta_k$, the norm of $\varepsilon$ from the quadratic subalgebra $K(\varepsilon)$ is $-d_k$, so $N(\varepsilon)=(-d_k)^3$, which is (H1). A root of $f_k$ in $K$ would be a root of $q_k$, which would make $d_k$ a square, or a root of $h_k$, which is irreducible over $\kv$ (Lemma~\ref{lem:fk}); this is (H2). The only monic quadratic factor of $f_k$ is $q_k$, because $c_kh_k$ is irreducible, and $\mu(T_{q_k})\neq0$ because its image in $H_{f_k}$ over $\kv$ is nontrivial; this is (H3).
\end{proof}

\subsection{The Abel-Prym map}
Let $L$ be a field containing $K$ and $P=(p,r,s)$ a point of $D_\delta(L)$, with $p\in L^3$. Bruin \cite[Lemma~6.1 and Proposition~6.2]{Bruin08} attaches to $P$ a point of the Jacobian of $Y^2=f_\delta(t)$ and gives the first polynomial of its Mumford representation \cite[Lemma~6.1(iii)]{Bruin08}: let $T$ be a point of the tangent line of $D_\delta$ at $P$ other than $P$, and let $a_i$ be the value at $T$ of the $i$-th quadric of~\eqref{eq:Ddelta}, $Q_i(T_p)-\delta\,(T_r^2,T_rT_s,T_s^2)_i$; then
\[
U_P(t)=t^2+2\frac{a_2}{a_3}t+\frac{a_1}{a_3}
\]
does not depend on~$T$. Bruin fixes the second polynomial by rationality. The construction of the second polynomial below, through a Pl\"ucker matrix, is ours. At a root $t_0$ of $U_P$, let $G_{t_0}$ be the $4\times4$ matrix $\mathrm{diag}(M_1+2t_0M_2+t_0^2M_3,-\delta)$ and let $\Omega=ab^{\mathsf T}-ba^{\mathsf T}$, for $a=(p,r+t_0s)$ and $b=(T_p,T_r+t_0T_s)$, be the Pl\"ucker matrix of a line in the quadric of $\mathbb P^3$ defined by $G_{t_0}$. Then $G_{t_0}\Omega G_{t_0}=Y_0\,{\star}\Omega$ for a scalar $Y_0$, where ${\star}$ is the Hodge dual, and the polynomial $V_P$ of degree at most one that takes the value $Y_0$ at each root $t_0$ satisfies $U_P\mid V_P^2-f_\delta$. We use the same construction for the data $(M_3,M_2,M_1)$ at $(p,s,r)$, which gives the reversed model: for $t=1/t'$, the matrix $M_1+2tM_2+t^2M_3$ is $t^2(M_3+2t'M_2+t'^2M_1)$.

A certificate at $P$ is a tangent vector $T$ as above, the proof that the tangent map of $D_\delta$ at $P$ has rank~$3$, the inequality $a_3\neq0$, a polynomial $V$ of degree at most one with $G_t\Omega G_t\equiv V\,{\star}\Omega$ modulo $U_P$ in $L[t]$, an entry of ${\star}\Omega$ invertible modulo $U_P$, and the divisibility $U_P\mid V^2-f_\delta$.

\begin{proposition}\label{prop:abel-prym}
Let $k\in\{0,1\}$. Every point $P\in D_{\delta_k}(K)$ has a certificate for the swapped data, and the class $\varphi_k(P)\in A_{f_k}(K)$ of the Mumford ideal $\langle U_P,Y-V\rangle$ does not depend on the certificate. The map $\varphi_k:D_{\delta_k}(K)\to A_{f_k}(K)$ satisfies $\varphi_k(\iota P)=-\varphi_k(P)$. If $\sigma:K\to L$ is a field embedding, the image under $\sigma$ of a certificate at $P$ is a certificate at $\sigma(P)$, and the classes correspond under $A_{f_k}(K)\to A_{\sigma f_k}(L)$. The values $\varphi_k(x_i)$ at the four known lifts are explicit Mumford ideals over~$K$.
\end{proposition}
\leanref[; the certificates are \path{cert_exists_fRev} in \path{FurioLombardo/Discharge/M3a/AbelPrymTotal.lean}, the independence of the certificate is \path{FurioLombardo.Discharge.M3a.AbelPrym.phiRev_eq_of_cert} in \path{FurioLombardo/Discharge/M3a/AbelPrym.lean}, the relation with $\iota$ is \path{FurioLombardo.Discharge.M3a.AbelPrym.phiRev_inv} in \path{FurioLombardo/Discharge/M3a/AbelPrym.lean}, the base change \path{phiRev_map} in \path{FurioLombardo/Discharge/M3a/CertMap.lean}, and the values at the known lifts are \path{phiRev_x0} to \path{phiRev_x3} in \path{FurioLombardo/Discharge/M3a/AbelPrymKnown.lean}]{FurioLombardo.Discharge.M3a.Bruin.phiK}{FurioLombardo/Discharge/M3a/ConcretePhi.lean}

\begin{proof}
Over any field $L$ of characteristic different from $2$, a certificate exists at a point $(p,r,s)$ of $D_\delta(L)$ with $(r,s)\neq(0,0)$ as soon as $f$ has no root in $L$ and its leading coefficient is not a square: the tangent map has rank $3$ because $(s^2M_1-2rsM_2+r^2M_3)p\neq0$, which follows from the absence of roots; a second tangent vector exists by counting dimensions; a nonzero minor of $(p,T_p)$ gives $a_3\neq0$ and the invertible entry; and $V$ is obtained by interpolation. For $f_k$ over $K$ these conditions hold by the proof of Proposition~\ref{prop:xT-kernel} and Lemma~\ref{lem:fk}, and $(r,s)\neq(0,0)$ because $Q_1,Q_2,Q_3$ have no common zero in $K^3$ other than $0$ (an explicit identity $x_j^4=\sum_ia_{j,i}Q_i$ with quadratic forms $a_{j,i}$ over $K$, for each coordinate $x_j$). Two certificates at the same point give the same $U_P$, since the tangent line is unique, and the same $V$, because the ruling identity and the invertible entry determine $V$ modulo~$U_P$. A certificate at $P$ gives one at $\iota P$ with $(U_P,-V)$, and the class of $\langle U,Y+V\rangle$ is the opposite of the class of $\langle U,Y-V\rangle$. The conditions of a certificate are polynomial identities, so they are preserved by field embeddings. The four values are computed exactly.
\end{proof}

We write $A_k=A_{f_k}$ and $T_k=T_{q_k}$. For $k=0$ the known lifts are $x_a=x_0$ and $x_b=x_2$, and for $k=1$ they are $x_a=x_1$ and $x_b=x_3$.

\section{Stoll's criterion at one place}\label{sec:stoll}

This section proves the abstract statement that turns global and local information on the group $A_k(K)$ into a constraint on the image of a global point; it is Stoll's Selmer group Chabauty \cite[Theorem~2.1]{Stoll17}, written at a single place and in terms of a logarithm.

Let $\mathsf A$, $\mathsf B$, $\mathsf V$ and $\mathsf M$ be abelian groups.

\begin{theorem}\label{thm:stoll-one-place}
Let $\iota:\mathsf A\to\mathsf B$, $\lambda:\mathsf B\to\mathsf V$ and $\pr:\mathsf V\to\mathsf M$ be homomorphisms, $T\in\mathsf A$, $\Gamma\subseteq\mathsf A$ a subgroup, $W\subseteq\mathsf V$ a subset and $x\in\mathsf A$. Assume:
\begin{enumerate}
\item[(loc)] if $\iota(Q)\in2\mathsf B$ for some $Q\in\mathsf A$, then $Q\in2\mathsf A$;
\item[(ker)] if $\lambda(b)=0$, then $b\in2\mathsf B$ or $b\in\iota(T)+2\mathsf B$;
\item[($\Gamma$)] $T\in\Gamma$, and $\pr\lambda\iota$ vanishes on $\Gamma$;
\item[(W)] for every $Q\in\mathsf A$, $\lambda\iota(Q)\in w+2\lambda(\mathsf B)$ for some $w\in W$;
\item[(sat)] if $Q\in\mathsf A$ and $\pr\lambda\iota(Q)\in2\mathsf M$, then $\lambda\iota(Q)\in\lambda\iota(\gamma)+2\lambda(\mathsf B)$ for some $\gamma\in\Gamma$;
\item[(M)] $\mathsf M$ has no $2$-torsion, and no nonzero element of $\mathsf M$ is divisible by every power of~$2$;
\item[(iii)] for all $n\ge0$, $z\in\mathsf M$ and $w\in W$ with $2^nz=\pr\lambda\iota(x)$ and $z-\pr(w)\in2\mathsf M$, one has $z\in2\mathsf M$.
\end{enumerate}
Then $\pr\lambda\iota(x)=0$.
\end{theorem}
\leanref{FurioLombardo.M4.stoll_log}{FurioLombardo/M4/Stoll.lean}

When $y=\pr\lambda\iota(x)$ is nonzero, (M) gives a largest $\nu$ with $y\in2^\nu\mathsf M$, and $y=2^\nu z$ with $z$ unique and $z\notin2\mathsf M$; condition (iii) then says that the class of $z$ in $\mathsf M/2\mathsf M$, the leading class of $y$, is not the class of an element of $\pr(W)$.

\begin{proof}
Suppose $y=\pr\lambda\iota(x)\neq0$, and let $\nu$ and $z$ be as above. We show by induction that for $0\le n\le\nu$ there are $t_n\in\Gamma$ and $Q_n\in\mathsf A$ with $x=t_n+2^nQ_n$. For $n=0$ take $t_0=0$ and $Q_0=x$. Let $n<\nu$. By ($\Gamma$), $y=2^n\pr\lambda\iota(Q_n)$, and since $\mathsf M$ has no $2$-torsion, $\pr\lambda\iota(Q_n)=2^{\nu-n}z\in2\mathsf M$. By (sat) there are $\gamma\in\Gamma$ and $b\in\mathsf B$ with $\lambda(\iota(Q_n-\gamma)-2b)=0$. By (ker), $\iota(Q_n-\gamma)$ or $\iota(Q_n-\gamma-T)$ lies in $2\mathsf B$, and by (loc) $Q_n=\gamma'+2Q'$ with $\gamma'=\gamma$ or $\gamma'=\gamma+T$, which lies in $\Gamma$ by ($\Gamma$). Then $x=(t_n+2^n\gamma')+2^{n+1}Q'$, which completes the induction. For $n=\nu$ we get $y=2^\nu\pr\lambda\iota(Q_\nu)$, so $\pr\lambda\iota(Q_\nu)=z$. By (W), $\lambda\iota(Q_\nu)=w+2\lambda(b)$ with $w\in W$, so $z-\pr(w)\in2\mathsf M$, and (iii) gives $z\in2\mathsf M$, a contradiction.
\end{proof}

In the application, $\mathsf A=A_k(K)$, $\mathsf B=A_k(\kv)$, $\iota$ is the map induced by $K\to\kv$, $\lambda$ is a logarithm on $A_k(\kv)$ with values in $\Zt^6$ (Section~\ref{sec:local}), $\mathsf M$ is the quotient of the image $\Lambda_k$ of $\lambda$ by the saturation $S_k$ of the span of the logarithms of the known points, and $x=\varphi_k(P)-\varphi_k(x_a)$ for a lift $P$ of a rational point. Condition (loc) comes from a bound on the image of the $x-T$ map (Section~\ref{sec:selmer}), (ker) from the torsion of $A_k(\kv)$ (Section~\ref{sec:local}), and (iii) from a covering of $C(\Qt)$ (Section~\ref{sec:cover}). The proof uses no generator of $A_k(K)$ and no bound on its rank: the induction stops at $\nu$, which is finite by Krull's intersection theorem for the finitely generated $\Zt$-module~$\mathsf M$.

\subsection{A form with a saturated subgroup}
In the application the subgroup $\Gamma$ is the preimage of a saturated submodule, and conditions (sat) and (M) follow from lattice data. Let $\Lambda$ be a finitely generated $\Zt$-module and $S\subseteq\Lambda$ a submodule. We say that $S$ is saturated in $\Lambda$ if $x\in\Lambda$ and $2x\in S$ imply $x\in S$. For $W\subseteq\Lambda$ and $y\in\Lambda$, condition (iii) at $y$ is the statement: for all $n\ge0$, $z\in\Lambda$ and $w\in W$ with $y-2^nz\in S$ and $z-w\in S+2\Lambda$, one has $z\in S+2\Lambda$. With $\mathsf M=\Lambda/S$ this is condition (iii) of Theorem~\ref{thm:stoll-one-place}.

\begin{lemma}\label{lem:leading-class}
Let $S$ be saturated in $\Lambda$ and $W\subseteq\Lambda$. Let $y=s_0+2^\nu z_0+2^{\nu+1}l$ with $s_0\in S$ and $z_0,l\in\Lambda$, where neither $z_0$ nor any $z_0-w$ with $w\in W$ lies in $S+2\Lambda$. Then $y\notin S$, and condition (iii) holds at~$y$.
\end{lemma}
\leanref[ and \path{FurioLombardo.M4.not_mem_of_leading}]{FurioLombardo.M4.condIII_of_leading}{FurioLombardo/M4/Leading.lean}

\begin{proof}
The quotient $\Lambda/S$ has no $2$-torsion because $S$ is saturated. In it, the image of $y$ is $2^\nu$ times the image $m$ of $z_0+2l$, and $m\notin2(\Lambda/S)$ because $z_0\notin S+2\Lambda$; in particular $y\notin S$. If $y-2^nz\in S$ with $n>\nu$, then $2^\nu(m-2^{n-\nu}\bar z)=0$, so $m=2^{n-\nu}\bar z$, which is impossible. If $n<\nu$, then $\bar z=2^{\nu-n}m$ and $z\in S+2\Lambda$. If $n=\nu$, then $\bar z=m$, so $z-z_0\in S+2\Lambda$, and $z-w\in S+2\Lambda$ would give $z_0-w\in S+2\Lambda$, which is excluded.
\end{proof}

\begin{corollary}\label{cor:stoll-sat}
Let $\iota:\mathsf A\to\mathsf B$ and $\lambda:\mathsf B\to\Zt^6$ be homomorphisms of abelian groups, $\Lambda=\lambda(\mathsf B)$ a $\Zt$-submodule, and $T\in\mathsf A$ with $2T=0$. Let $\gamma_a,\gamma_b\in\mathsf A$, $a=\lambda\iota(\gamma_a)$ and $b=\lambda\iota(\gamma_b)$, let $S$ be a saturated submodule of $\Lambda$ containing $a$ and $b$, and let $W\subseteq\Lambda$. Assume (loc) and (ker) of Theorem~\ref{thm:stoll-one-place}, that every $\lambda\iota(Q)$ with $Q\in\mathsf A$ lies in $w+2\Lambda$ for some $w\in W$, and that there are integers $n,m$ and $r\in\Lambda$ with $4r=na+mb$ and $S\subseteq\Zt a+\Zt b+\Zt r$. If $x\in\mathsf A$ and condition (iii) holds at $\lambda\iota(x)$, then $\lambda\iota(x)\in S$.
\end{corollary}
\leanref[, in the form used below, for all the points of one twist at once; it applies \path{FurioLombardo.M3a.stoll_sat} in the same file, which is Theorem~\ref{thm:stoll-one-place} with $\mathsf M=\Lambda/S$]{FurioLombardo.M3a.twist_sat}{FurioLombardo/M3a/StollSat.lean}

\begin{proof}
Since $\Zt^6$ has no $2$-torsion, $\lambda\iota(T)=0$. We first show that $r=\lambda\iota(R)$ for some $R\in\mathsf A$. If $\gamma\in\mathsf A$ and $\lambda\iota(\gamma)=2\lambda(b')$ with $b'\in\mathsf B$, then by (ker) $\iota(\gamma)-2b'$ lies in $2\mathsf B$ or in $\iota(T)+2\mathsf B$, so by (loc) $\gamma$ or $\gamma-T$ is $2R'$ with $R'\in\mathsf A$, and then $2\lambda\iota(R')=2\lambda(b')$ gives $\lambda\iota(R')=\lambda(b')$. Applied to $\gamma=n\gamma_a+m\gamma_b$, whose image is $4r$, and then to the resulting point, whose image is $2r$, this gives $R$. Since $4r\in S$ and $S$ is saturated, $r\in S$.

Let $\Gamma\subseteq\mathsf A$ be the preimage of $S$ under $\lambda\iota$. It contains $T$, $\gamma_a$, $\gamma_b$ and $R$, and every element of $S$ is $\alpha a+\beta b+\rho r$ with $\alpha,\beta,\rho\in\Zt$; writing each coefficient as an integer plus twice an element of $\Zt$ shows that every element of $S$ is $\lambda\iota(\gamma)+2s$ with $\gamma\in\Gamma$ and $s\in S$. Now apply Theorem~\ref{thm:stoll-one-place} with $\lambda$ viewed as a map to $\mathsf V=\Lambda$, with $\mathsf M=\Lambda/S$ and with $\pr$ the projection. Conditions ($\Gamma$) and (W) hold by construction. For (sat), if $\lambda\iota(Q)=s+2l$ with $s\in S$ and $l\in\Lambda$, then $\lambda\iota(Q)=\lambda\iota(\gamma)+2(s'+l)$ with $\gamma\in\Gamma$ and $s'+l\in\Lambda=\lambda(\mathsf B)$. For (M), $\Lambda/S$ has no $2$-torsion because $S$ is saturated, and no nonzero element divisible by every power of $2$ by Krull's intersection theorem for the finitely generated $\Zt$-module $\Lambda/S$. So $\pr\lambda\iota(x)=0$.
\end{proof}

\section{The Selmer bound}\label{sec:selmer}

This section proves that the image of $\mu$ on $A_k(K)$ is contained in a group generated by four explicit classes whose images at $v$ are known (Theorem~\ref{thm:selmer-basis}). This gives conditions (loc) and (W) of Corollary~\ref{cor:stoll-sat}. The proof is a $2$-descent in the manner of Schaefer \cite{Schaefer95} and of Poonen and Schaefer \cite{PS97}: a global bound from $S$-units, then local conditions at finitely many places, each given by a finite certificate.

Fix $k\in\{0,1\}$ and write $f=f_k$, $q=q_k$, $h=h_k$ and $H=H_{f}$. The algebra $\mathcal A_f=K[X]/(f)$ is isomorphic to $L\times N$, where $L$ is the quadratic extension of $K$ generated by a root $\alpha$ of $q$, and $N$ is the quadratic extension of $L$ generated by a root $\beta'$ of one of the two quadratic factors of $h$ over $L$ (the factor $\Psi$ of the proof of Lemma~\ref{lem:fk}); the isomorphism sends $\Theta$ to $(\alpha,\beta')$. The fields $L$ and $N$ have degrees $42$ and $84$ over~$\Q$.

\subsection{The global bound}

\begin{lemma}\label{lem:schaefer-local}
Let $M$ be a field with a discrete valuation $v$, valuation ring $R$ and uniformizer $\varpi$. Let $F\in R[X]$ have degree~$6$ and a unit leading coefficient, and let $\theta\in M$ be a root of $F$ with $F'(\theta)$ a unit. Let $u,w,z\in M[X]$ with $\deg u=2$, $\deg w\le1$, $w^2-F=uz$ and $u(\theta)\neq0$. Let $c\in M^\times$ be such that $c^{-1}u$ has coefficients in $R$, not all in the maximal ideal. Then $v(u(\theta)/c)$ is even.
\end{lemma}
\leanref[; the form used in the proof of Proposition~\ref{prop:schaefer-global}, with the parity as a square class, is \path{FurioLombardo.Discharge.SelmerBasis.schaefer_local_sq} in the same file]{FurioLombardo.Discharge.SelmerBasis.schaefer_local}{FurioLombardo/Discharge/SelmerBasis/SchaeferLocal.lean}

\begin{proof}
The root $\theta$ is integral, and $F=(X-\theta)g$ with $g\in R[X]$ and $g(\theta)=F'(\theta)$ a unit, so $(X-\bar\theta)^2$ does not divide $\bar F$; bars denote reduction. Put $U=c^{-1}u$ and $Z=cz$, so that $w^2-F=UZ$, and write $Z=\varpi^tZ'$ with $Z'$ primitive. Comparing contents, and using that the reduction of a product of primitive polynomials is nonzero, gives three cases. If $w$ is $0$ or has coefficients in the maximal ideal, then $t=0$ and $\bar U\bar Z'=-\bar F$. If $w$ is primitive, then $t=0$ and $\bar U\bar Z'=\bar w^2-\bar F$, since the leading coefficient of $F$ is a unit. If $w=\varpi^{s}w'$ with $s<0$ and $w'$ primitive, then $t=2s$ and $\bar U\bar Z'=\bar w'^2$. In every case $t$ is even.

Evaluating at $\theta$ gives $U(\theta)\varpi^tZ'(\theta)=w(\theta)^2$, so $v(U(\theta))$ is even when $Z'(\theta)$ is a unit. Next, $\Res(U,F)=U(\theta)\Res(U,g)$, and $\Res(U,F)=\Res(U,w^2-UZ)$ is a nonzero square times a power of the leading coefficient of $U$ with even exponent, so $v(U(\theta))$ is even when $\Res(U,g)$ is a unit, that is when $\bar U$ and $\bar g$ are coprime. If $U(\theta)$ is a unit there is nothing to prove. Otherwise, if $Z'(\theta)$ is not a unit either, then $(X-\bar\theta)^2$ divides $\bar U\bar Z'$. In the first case this contradicts $(X-\bar\theta)^2\nmid\bar F$. In the second, $X-\bar\theta$ divides $\bar F$ and $\bar w^2-\bar F$, hence $\bar w$, so $(X-\bar\theta)^2$ divides $\bar w^2$ and $\bar F$, a contradiction. In the third, $\bar U\bar Z'=\bar w'^2$ has degree at most $2$ and both factors vanish at $\bar\theta$, so $\bar U=e(X-\bar\theta)$ with $e\neq0$, which is prime to $\bar g$ because $\bar g(\bar\theta)\neq0$.
\end{proof}

The lemma uses no completion and no hypothesis on the residue characteristic.

\begin{proposition}\label{prop:schaefer-global}
Let $\langle u,Y-w\rangle$ be a Mumford ideal for $f$ over $K$ with $u$ prime to~$f$. There is $c\in K^\times$ with the following property: for every number field $E$ containing $K$, every root $\theta\in E$ of $f$ and every prime $\mathfrak P$ of $E$ not above $2$ or $7$, the valuation of $u(\theta)/c$ at $\mathfrak P$ is even.
\end{proposition}
\leanref[; the abstract form, with the certificates as hypotheses, is \path{schaefer_global_of_certs} in \path{FurioLombardo/Discharge/SelmerBasis/SchaeferModel.lean}]{FurioLombardo.Discharge.SelmerBasis.schaeferGlobal}{FurioLombardo/Discharge/SelmerBasis/SchaeferGlobal.lean}

\begin{proof}
The coefficients of $f$ are integral outside $14$, but $f$ does not satisfy the hypotheses of Lemma~\ref{lem:schaefer-local} at every prime outside $14$: modulo $\mathfrak q_3$ and $\mathfrak q_{439}$ its six roots coincide, and at the five primes above $23$ its leading coefficient is not a unit. So we change the model. Let $\varpi\in\OK$ generate $\mathfrak q_3\mathfrak q_{439}$, of norm $1317$, let $a_T=-469$, and put
\[
f_1(X)=\varpi^{-6}f(a_T+\varpi X),\qquad U_1(X)=\varpi^{-2}u(a_T+\varpi X),
\]
and $w_1=\varpi^{-3}w(a_T+\varpi X)$, so that $U_1$ is monic, $U_1\mid w_1^2-f_1$, and $u(\theta)=\varpi^2U_1(\theta_1)$ with $\theta_1=(\theta-a_T)/\varpi$. Let $c$ generate the fractional ideal generated by $1$ and the coefficients of $U_1$; it exists because $\OK$ is principal (Proposition~\ref{prop:classnumber}). Let $\mathfrak P$ be a prime of $E$ not above $2$ or~$7$. If the leading coefficient $f_6$ of $f$ is not in $\mathfrak P$, then $f_1$ is $\mathfrak P$-integral with unit leading coefficient $f_6$, and $f_1'(\theta_1)$ is a unit by an explicit identity $A_1f_1+B_1f_1'=e_1f_6$ with $A_1,B_1$ integral outside $14$ and $e_1$ dividing a power of~$14$. Lemma~\ref{lem:schaefer-local} for $(f_1,U_1,\theta_1,c)$ gives the parity for $U_1(\theta_1)/c$, hence for $u(\theta)/c$. If $f_6\in\mathfrak P$, one of the reversed models $f_2=X^6f_1(i+1/X)$ with $i\in\{0,1,2\}$ and $U_1(i)\neq0$ has unit leading coefficient $f_1(i)$ and a unit derivative at $\theta_2=1/(\theta_1-i)$, by an explicit identity $yf_6+zf_1(i)r_2=14^N$ and a B\'ezout identity for $f_2$ with constant $r_2$. Lemma~\ref{lem:schaefer-local} applies to $f_2$, to the monic polynomial $U_2=U_1(i)^{-1}X^2U_1(i+1/X)$, to $w_2=X^3w_1(i+1/X)$ reduced modulo $U_2$ to degree at most one, and to the constant $c/U_1(i)$, and the identity $U_1(\theta_1)=(\theta_1-i)^2U_1(i)U_2(\theta_2)$ gives the parity. Every input is an identity over $K$ checked by exact computation; the proof splits on whether $f_6\in\mathfrak P$ and never names the primes involved.
\end{proof}

\begin{lemma}\label{lem:sunit-span}
The class groups of $L$ and $N$ have odd order. Let $S_L$ and $S_N$ be the sets of primes of $L$ and $N$ above $2$ and $7$. Every element of $L^\times$ with even valuation at every prime outside $S_L$ is a product of $29$ explicit elements of $L$ times a square, and every element of $N^\times$ with even valuation at every prime outside $S_N$ is a product of $53$ explicit elements of $N$ times a square.
\end{lemma}
\leanref[ and \path{sUnitSpan_N84}; the class groups are \path{clTwo_L42} and \path{clTwo_N84} in the same file]{FurioLombardo.Discharge.SelmerBasis.sUnitSpan_L42}{FurioLombardo/Discharge/SelmerBasis/SUnitSpan.lean}

\begin{proof}
Let $L'/K'$ be a quadratic extension of number fields with $\Gal(L'/K')=\langle\tau\rangle$. Chevalley's ambiguous class number formula \cite[p.~406]{Chevalley33} gives
\[
\#\Cl(L')^{\langle\tau\rangle}=\frac{h_{K'}\,2^{t-1}}{[E_{K'}:E_{K'}\cap N_{L'/K'}(L'^\times)]},
\]
where $E_{K'}$ is the unit group of $K'$ and $t$ is the number of places of $K'$, finite or real, that ramify in~$L'$. The index is a power of $2$, because $E_{K'}^2$ consists of norms. So if $h_{K'}$ is odd and the index is at least $2^{t-1}$, the group $\Cl(L')^{\langle\tau\rangle}$ has odd order, and the $2$-part of $\Cl(L')$ has no nonzero fixed point under $\tau$. Since $\tau$ has order $2$, the number of fixed points in the $2$-part is congruent to its order modulo $2$, so the $2$-part is trivial.

We have $L=K(\sqrt d)$ and $N=L(\sqrt e)$ for explicit $d\in K$ and $e\in L$. For $L/K$, with $h_K=1$: only the places $v$ and $w_6$ ramify, since $d$ is positive at the three real places, so $t\le2$, and an explicit unit of $K$ is not a norm from $L$, which is decided by its residue modulo $4$ at~$v$. For $N/L$, with $\Cl(L)$ now of odd order: one prime above $7$ and three real places of $L$ ramify, so $t\le4$, and three explicit units of $K$ have the diagonal sign pattern at these three real places, so they are independent modulo norms and the index is at least $2^3$.

As in the proof of Lemma~\ref{lem:sunits-K}, an element of $L^\times$ with even valuations outside $S_L$ is then an $S_L$-unit times a square. The unit group of $L$ has rank at most $23$ and $S_L$ has at most five primes, so the $S_L$-units modulo squares form a space of dimension at most $1+23+5=29$. The $29$ listed elements lie in it and are independent modulo squares, by their quadratic residue characters at auxiliary primes, so they span it. For $N$ the same argument, with rank at most $44$ and at most eight primes, gives $1+44+8=53$.
\end{proof}

We write $g_1,\dots,g_{82}$ for the classes in $H$ of the elements $(\gamma,1)$ and $(1,\gamma')$ of $L\times N$, with $\gamma$ one of the $29$ elements of $L$ and $\gamma'$ one of the $53$ elements of~$N$.

\begin{lemma}\label{lem:global-span}
For every $Q\in A_k(K)$, $\mu(Q)$ lies in the subgroup of $H$ generated by $g_1,\dots,g_{82}$.
\end{lemma}
\leanref{FurioLombardo.Discharge.SelmerBasis.GlobalK21.globalSpan_gK}{FurioLombardo/Discharge/SelmerBasis/GlobalSpan.lean}

\begin{proof}
The classes $0$ and $T_k$ are handled directly. Every other class of $A_k(K)$ is the class of a Mumford ideal $\langle u,Y-w\rangle$ with $u$ prime to $f$, as in the proof of Proposition~\ref{prop:xT-kernel}, and $\mu(Q)$ is the class of $u(\Theta)$, which corresponds to $(u(\alpha),u(\beta'))$. By Proposition~\ref{prop:schaefer-global}, applied with $E=L$ and $E=N$, the elements $u(\alpha)/c$ and $u(\beta')/c$ have even valuations outside $S_L$ and $S_N$, so by Lemma~\ref{lem:sunit-span} each is a product of the listed elements times a square. Since $c\in K^\times$, the class of $u(\Theta)$ in $H$ is the class of $(u(\alpha)/c,u(\beta')/c)$.
\end{proof}

\subsection{The local conditions}

For a place $w$ of $K$ let $K_w$ be the completion, $H_w$ the group $H_{f}$ for $f$ over $K_w$, and $\mu_w$ the $x-T$ map on $A_{f}(K_w)$; the map $K\to K_w$ induces $H\to H_w$, compatible with $\mu$ and $\mu_w$. We use the places
\[
\Sigma_0=\{v,w_{12},w_6,\rho_1,\rho_2\},\qquad\Sigma_1=\{v,w_{12},w_6,w_7\},
\]
for $k=0$ and $k=1$, where $\rho_1$ and $\rho_2$ are two of the three real places of~$K$.

\begin{lemma}\label{lem:count}
Let $w$ be one of $v$, $w_{12}$, $w_6$, and also $w_7$ when $k=1$, and let $n_w=7$, $26$, $14$, $3$ respectively. Then $\mu_w$ takes at most $2^{n_w}$ values on $A_f(K_w)$.
\end{lemma}
\leanref[; the other places are \path{countBound_w2}, \path{countBound_w3} and \path{countBound_w7} in \path{AtW2.lean}, \path{AtW3.lean} and \path{At7.lean} of the same directory]{FurioLombardo.Discharge.SelmerBasis.Count.countBound_v}{FurioLombardo/Discharge/SelmerBasis/Count/AtV.lean}

\begin{proof}
Since $H_w$ has exponent $2$, $\mu_w$ factors through $A_f(K_w)/2A_f(K_w)$. Let $\mathcal O_w$ be the ring of integers of $K_w$ and $e_w$ its ramification index over $\Q_p$. The formal group of the Jacobian at a point of the curve over $K_w$ gives a subgroup of finite index in $A_f(K_w)$ isomorphic to $\mathcal O_w^2$, which has no $2$-torsion; the construction is the one of Proposition~\ref{prop:log} below, carried out over any complete field. Hence $[A_f(K_w):2A_f(K_w)]=[\mathcal O_w^2:2\mathcal O_w^2]\cdot\#A_f(K_w)[2]$. The first factor is $2^{2e_w}$ at the places above $2$, whose residue fields are $\F_2$, and $1$ at~$w_7$. A point of order dividing $2$ is $0$ or $T_u$ for a monic quadratic factor $u$ of $f$ over $K_w$. At $v$ the only such factor is $q$ (Lemma~\ref{lem:fk}), so $\#A_f(K_v)[2]\le2$. At $w_{12}$ and $w_6$, $d_k$ is not a square, so $f$ has no root in $K_w$ and has at most three monic quadratic factors, which gives at most $4$. At $w_7$ the norm of the discriminant of the factor $\Psi$ of $h$ in the proof of Lemma~\ref{lem:fk} is not a square, so $f$ has an irreducible quadratic factor and at most seven monic quadratic factors, which gives at most $8$. The bounds are $2^{6+1}$, $2^{24+2}$, $2^{12+2}$ and $2^{0+3}$.
\end{proof}

\begin{lemma}\label{lem:local-images}
Let $w\in\Sigma_k$. There is an explicit matrix $C_w$ over $\F_2$ with $82$ columns such that, for every $a\in\Z^{82}$, if the image of $\prod_sg_s^{a_s}$ in $H_w$ lies in $\mu_w(A_f(K_w))$, then $C_wa\equiv0\pmod2$. The matrices have $11$, $39$, $21$ and $5$ rows at $v$, $w_{12}$, $w_6$ and $w_7$, and $2$ rows at each real place.
\end{lemma}
\par\smallskip{\raggedright\noindent\textit{Lean:} in \path{FurioLombardo/Discharge/SelmerBasis}, namespace \path{FurioLombardo.Discharge.SelmerBasis}: at $v$ \path{V.coordCond_v} in \path{VCoord.lean} with \path{V.imageIn_v} in \path{VPlace.lean}; at $w_{12}$ \path{W2.coordCond_w2} in \path{W2Place.lean}; at $w_6$ \path{W3.coordCond_w3} with \path{W3.imageIn_w3} in \path{W3Place.lean}; at $w_7$ \path{W7.coordCond_w7} in \path{W7Place.lean}; at the real places \path{RealRoots.coordCond_place} in \path{RealCoord.lean}, with the signs \path{RealRoots.signOK} in \path{RealSigns.lean} and the image \path{RealRoots.imageIn_place} in \path{RealInst.lean}.\par}\smallskip

\begin{proof}
Let $w$ be a finite place. Over $K_w$, $f$ is its leading coefficient times a product of irreducible factors $g_j$, of degrees $2,4$ at $v$, $2,2,2$ at $w_{12}$ and $w_6$, and $1,1,1,2,1$ at $w_7$. Let $F_j=K_w[X]/(g_j)$ and $G=\prod_jF_j^\times$. Evaluation at roots of the $g_j$, given by Hensel's lemma from explicit approximations, is a ring homomorphism $\mathcal A_f\otimes K_w\to\prod_jF_j$; only this is used, and it maps $K_w^\times(\mathcal A_f\otimes K_w)^{\times2}$ into $K_w^\times G^2$. Square classes in $G$ are read from valuations and from the depth of units: in a field $F_j$ with ramification index $e_j$ over $\Qt$ and uniformizer $\Pi$, a unit $1+\Pi^td$ with $d$ a unit and $t$ odd, $t<2e_j$, is not a square, a unit $1+4m$ is a square only if the residue of $m$ has trace $0$ to $\F_2$, and a unit $\equiv1\pmod{4\Pi}$ is a square; at $w_7$ a unit is a square exactly when its residue is. From this, explicit certificates give each of the following elements coordinates over $\F_2$ in a common space: generators of $K_w^\times$ modulo squares, the images of $n_w$ explicit points of $A_f(K_w)$, and the images of $g_1,\dots,g_{82}$. The images of the $n_w$ points are independent modulo the image of $K_w^\times$, so by Lemma~\ref{lem:count} they span $\mu_w(A_f(K_w))$. The rows of $C_w$ are linear forms that vanish on the coordinates of $K_w^\times$ and of the $n_w$ points, evaluated on the coordinates of the~$g_s$.

Let $\rho$ be a real place, among the two that we use. The leading coefficient of $f$ is negative at $\rho$, and $f$ has exactly four real roots $r_1<r_2<r_3<r_4$, located by sign changes at rational points; so $f>0$ exactly on $(r_1,r_2)\cup(r_3,r_4)$. The algebra $\mathcal A_f\otimes\mathbb R$ is $\mathbb R^4\times\mathbb C$, and a class in $H_\rho$ is determined by its sign vector at $r_1,\dots,r_4$ modulo the diagonal. Let $\langle u,Y-w\rangle$ be a Mumford ideal with $u$ prime to $f$. If $u$ has no real root, all its signs are positive. If $u$ has real roots $x_1,x_2$, then $f(x_j)=w(x_j)^2>0$, so each $x_j$ lies in $(r_1,r_2)$ or $(r_3,r_4)$, and the sign vector of $u(r_i)=(r_i-x_1)(r_i-x_2)$ is $(+,+,+,+)$ or $(+,-,-,+)$. The classes $T_u$, and the ideals where $u$ shares one root with $f$, reduce to these cases. So the image of $\mu_\rho$ is the group of classes $\eta$ with $\eta(r_1)\eta(r_4)>0$ and $\eta(r_2)\eta(r_3)>0$, and the two rows of $C_\rho$ are the two sign conditions, read from the signs of the $82$ elements at the real roots, which are certified by interval arithmetic.
\end{proof}

\begin{theorem}\label{thm:selmer-basis}
For $k=0,1$ there are classes $\gamma_1,\dots,\gamma_4\in H$ such that $\mu(A_k(K))$ is contained in the subgroup they generate, and such that the image of $\gamma_j$ in $H_v$ is $\prod_{i=1}^7\mu_v(D_i)^{B_{ij}}$, where $D_1,\dots,D_7$ are the points of Lemma~\ref{lem:local-divisors} and $B$ is an explicit $7\times4$ integer matrix with a left inverse modulo~$2$.
\end{theorem}
\leanref[ and \path{selmerBasis_1}; the left inverse of $B$ modulo $2$ is \path{FurioLombardo.M4.T0.hSBL} and \path{FurioLombardo.M4.T1.hSBL} in \path{FurioLombardo/M4/Checks.lean}]{FurioLombardo.Discharge.SelmerBasis.Assembly.selmerBasis_0}{FurioLombardo/Discharge/SelmerBasis/Assembly.lean}

\begin{proof}
Let $Q\in A_k(K)$. By Lemma~\ref{lem:global-span}, $\mu(Q)=\prod_sg_s^{a_s}$ for some $a\in\Z^{82}$. Its image in $H_w$ is $\mu_w(Q)$, so $C_wa\equiv0\pmod2$ for every $w\in\Sigma_k$ by Lemma~\ref{lem:local-images}. The stacked matrix of the $C_w$ has $75$ rows for $k=0$ and $76$ rows for $k=1$, and rank $63$ in both cases. An explicit certificate shows that its kernel is spanned by four vectors $b_1,\dots,b_4$ and fifteen vectors $\kappa_1,\dots,\kappa_{15}$ such that $\prod_sg_s^{(\kappa_l)_s}=1$ in $H$; these products are classes of elements of $K^\times$ times squares, which is checked exactly. So $\mu(Q)$ lies in the group generated by the $\gamma_j=\prod_sg_s^{(b_j)_s}$. The images of the $\gamma_j$ in $H_v$ are computed in the coordinates of the proof of Lemma~\ref{lem:local-images} at~$v$ and written in the basis $\mu_v(D_1),\dots,\mu_v(D_7)$.
\end{proof}

\begin{corollary}\label{cor:selmer-consequences}
For $k=0,1$:
\begin{enumerate}
\item if $Q\in A_k(K)$ and its image in $A_k(\kv)$ lies in $2A_k(\kv)$, then $Q\in2A_k(K)$;
\item for every $Q\in A_k(K)$ there are $\epsilon\in\{0,1\}^4$ and $b\in A_k(\kv)$ such that the image of $Q$ in $A_k(\kv)$ is $\sum_i(B\epsilon)_iD_i+2b$.
\end{enumerate}
\end{corollary}
\leanref[, with the hypotheses given by Proposition~\ref{prop:xT-kernel} and by \path{FurioLombardo.Discharge.M3a.selmerInjective_of_indep} in \path{FurioLombardo/Discharge/M3a/SelmerK2.lean}; (2) is \path{FurioLombardo.Discharge.M4Box.hSel_of_selmerBasis} in \path{FurioLombardo/Discharge/M4Box/SelLoc.lean}]{FurioLombardo.M3a.Route.hloc_jac}{FurioLombardo/M3a/Route.lean}

\begin{proof}
Write $\mu(Q)=\prod_j\gamma_j^{\epsilon_j}$ with $\epsilon\in\{0,1\}^4$, by Theorem~\ref{thm:selmer-basis}. The image of $Q$ in $A_k(\kv)$ has $x-T$ image $\prod_i\mu_v(D_i)^{(B\epsilon)_i}$. For (1), this image is trivial, so $B\epsilon\equiv0\pmod2$ by the independence of the $\mu_v(D_i)$ (Lemma~\ref{lem:local-divisors}), hence $\epsilon=0$ by the left inverse of $B$. Then $\mu(Q)=1$, and $Q\in2A_k(K)$ by Proposition~\ref{prop:xT-kernel}. For (2), the image of $Q$ minus $\sum_i(B\epsilon)_iD_i$ lies in the kernel of $\mu_v$. This kernel is $2A_k(\kv)$: the group $A_k(\kv)/2A_k(\kv)$ has at most $2^7$ elements by Lemma~\ref{lem:count}, and $\mu_v$ takes at least $2^7$ values on it, the products of the independent classes $\mu_v(D_i)$.
\end{proof}

\section{The local group at \texorpdfstring{$v$}{v}}\label{sec:local}

This section describes $A_k(\kv)$ through a logarithm with values in $\Zt^6$ and computes the lattice data that enter Corollary~\ref{cor:stoll-sat}.

We use for $\kv$ the field $\Qt[x]/(E)$, where $E\in\Z[x]$ is an explicit Eisenstein cubic with $E\equiv x^3+2\pmod4$, together with the embedding $K\to\kv$ that sends $\beta$ to the root of $g$ given by Hensel's lemma at an explicit integral approximation $\theta_0$, with $|\beta-\theta_0|\le2^{-33}$. The ring of integers $\Ov$ has the basis $1,\pi,\pi^2$ over $\Zt$, where $\pi$ is the class of~$x$, and we identify $\Ov^2$ with $\Zt^6$ through this basis. This field is the completion of $K$ at $v$; the proofs only use that it is a complete field with an embedding of~$K$.

\begin{lemma}\label{lem:local-2-torsion}
Every point of $A_k(\kv)$ whose order is a power of $2$ is $0$ or $T_k$.
\end{lemma}
\leanref[, with the inputs \path{FurioLombardo.Discharge.M3a.Bruin.l1Inputs} and the local facts \path{localFacts_Kv} in \path{FurioLombardo/Discharge/M3a/LocalFacts.lean}; it is applied in \path{FurioLombardo.Discharge.M4Log.two_torsion_Kv}]{FurioLombardo.Discharge.M3a.localTwoTorsion_of_L1Inputs}{FurioLombardo/Discharge/M3a/Refined.lean}

\begin{proof}
A point of order dividing $2$ is $0$ or $T_u$ for a monic quadratic factor $u$ of $f_k$, and over $\kv$ the only one is $q_k$, because $c_kh_k$ is irreducible over $\kv$ (Lemma~\ref{lem:fk}). The image $\mu_v(T_k)$ is the class of $(q_k-c_kh_k)(\Theta)$, and an explicit square class computation shows that it is not of the form $c'y^2$ with $c'\in\kv^\times$. A point of order $4$ would give $T_k\in2A_k(\kv)$, hence $\mu_v(T_k)=1$, which is not the case.
\end{proof}

\begin{proposition}\label{prop:log}
For $k=0,1$ there are a homomorphism $\lambda_k:A_k(\kv)\to\Zt^6$, a subgroup $B_1$ of finite index in $A_k(\kv)$ and an integer $N$ such that $\lambda_k$ is injective on $B_1$ and $\lambda_k(B_1)$ contains $2^N\Zt^6$.
\end{proposition}
\leanref[; the properties are \path{FurioLombardo.Discharge.M4Box.DCertData.chartLog} in \path{FurioLombardo/Discharge/M4Box/DCertData/Values.lean}]{FurioLombardo.Discharge.M4Log.lam}{FurioLombardo/Discharge/M4Log/Defs.lean}

\begin{proof}[Construction]
Let $a_k=k$ and let $b_k\in\kv$ with $b_k^2=f_k(a_k)$. For $t=(t_0,t_1)\in\kv^2$ close to $0$, the Mumford ideal $\langle(X-a_k)^2+t_0(X-a_k)+t_1,\ Y-V_t\rangle$, with $V_t$ the square root of $f_k$ modulo the first generator that is close to the tangent line at $(a_k,b_k)$, defines a class near twice the class of the point $(a_k,b_k)$; let $\psi(z)$ be its class minus the class $E_0$ at $t=0$, for $t=\pi^{M}z$. For $M=8$ ($k=0$) or $M=0$ ($k=1$), the group law read through $\psi$ is a formal group law $\mathcal G$ in two variables with coefficients in $\Ov$, and $\psi$ is an injective homomorphism from the group $\mathcal G(\pi^4\Ov)$ onto a subgroup $B_1$ of finite index. On $B_1$ put $\lambda_k(\psi(z))=\pi^{M}\mathrm A\log_{\mathcal G}(z)$, where $\log_{\mathcal G}$ is the formal logarithm and $\mathrm A$ is the explicit $2\times2$ matrix that expresses, in the coordinates $t$, the linear part of the logarithm attached to the differentials $dx/y$ and $x\,dx/y$ of $y^2=f_{\delta_k}(x)$. Since $B_1$ has finite index, $\lambda_k$ extends uniquely to $A_k(\kv)$ with values in $\kv^2$, and it is bounded. Its values at the seven points $D_i$ of Lemma~\ref{lem:local-divisors} are integral by the certificates of Proposition~\ref{prop:lattice}; since the $D_i$ generate $A_k(\kv)$ modulo $2A_k(\kv)$, iterating $b=\sum c_iD_i+2b'$ shows that every value is integral. Through the basis $1,\pi,\pi^2$ this gives $\lambda_k$ with values in $\Zt^6$.
\end{proof}

We do not identify $\lambda_k$ with the analytic logarithm of the Jacobian; the proof uses only Proposition~\ref{prop:log} and the values of $\lambda_k$ computed from the chart.

\begin{lemma}\label{lem:local-divisors}
For $k=0,1$ there are seven explicit points $D_1,\dots,D_7$ of $A_k(\kv)$ whose images under $\mu_v$ are independent in $H_v$. Consequently the $D_i$ generate $A_k(\kv)/2A_k(\kv)$, condition (ker) of Theorem~\ref{thm:stoll-one-place} holds for $\lambda_k$ and $T_k$, and $\Lambda_k=\lambda_k(A_k(\kv))$ is the $\Zt$-module spanned by $\lambda_k(D_1),\dots,\lambda_k(D_7)$.
\end{lemma}
\leanref[; the generation modulo $2$ is \path{FurioLombardo.Discharge.M4Log.genModTwo_Dpt} in \path{FurioLombardo/Discharge/M4Log/LamInt.lean}, condition (ker) is \path{FurioLombardo.Discharge.Analytic.hKer_of_chart} in \path{FurioLombardo/Discharge/Analytic/LogChart.lean}, and the description of $\Lambda_k$ is \path{FurioLombardo.Discharge.M4Box.hLam_k} in \path{FurioLombardo/Discharge/M4Box/Final.lean}]{FurioLombardo.Discharge.SelmerSpan.indep_Dpt}{FurioLombardo/Discharge/SelmerSpan/Indep.lean}

\begin{proof}
Each $D_i$ is a Mumford ideal $\langle X^2+pX+r,\ Y-V\rangle$ with $p,r\in K$ exact and $V$ built from two square roots in $\kv$ given by Hensel's lemma. For each of the $127$ nonzero combinations of the $\mu_v(D_i)$, a certificate shows that the product is not of the form $c'y^2$: a norm of it to $\kv$, or to one of two extensions of $\kv$ given by explicit square roots, is not a square, which is decided at finite precision. The group $A_k(\kv)/2A_k(\kv)$ has at most $2^7$ elements by Lemma~\ref{lem:count} (in the form $[A_k(\kv):2A_k(\kv)]=[\Zt^6:2\Zt^6]\cdot2$ through $B_1$), so the seven points generate it. A point killed by $\lambda_k$ has a multiple in $B_1$ killed by $\lambda_k$, hence is torsion. A torsion point is the sum of a point of odd order, which lies in $2A_k(\kv)$, and a point of $2$-power order, which is $0$ or $T_k$ by Lemma~\ref{lem:local-2-torsion}; this is (ker). Finally $\Lambda_k$ is a finitely generated $\Zt$-module spanned by the $\lambda_k(D_i)$ modulo $2\Lambda_k$, and Nakayama's lemma gives the last claim.
\end{proof}

\begin{proposition}\label{prop:lattice}
For $k=0,1$ let $a=\lambda_k(\varphi_k(x_a))$ and $b=\lambda_k(\varphi_k(x_b))$, the images being taken in $A_k(\kv)$, and let $S_k$ be the saturation of $\Zt a+\Zt b$ in $\Lambda_k$. Let $\ell_1,\dots,\ell_7$ be integer vectors with $\lambda_k(D_i)-\ell_i\in8\Zt^6$, fixed in the proof below, and let $W_k$ be the set of the $16$ integer vectors $\sum_i(B\epsilon)_i\ell_i$ with $\epsilon\in\{0,1\}^4$. Then:
\begin{enumerate}
\item $4\Zt^6\subseteq\Lambda_k$, and $\Lambda_k=\{x\in\Zt^6:G_kx\in4\Zt^6\}$ for an explicit integer matrix~$G_k$;
\item there are integers $n,m$ and $r\in\Lambda_k$ with $4r=na+mb$ and $S_k\subseteq\Zt a+\Zt b+\Zt r$;
\item there are an explicit $4\times6$ integer matrix $P_k$ and an integer $r_k$ ($r_0=41$, $r_1=46$) such that, for $x\in\Lambda_k$ and $n\le r_k$, $x\in S_k+2^n\Lambda_k$ if and only if $2^{n+2}$ divides every coordinate of $P_kx$.
\end{enumerate}
\end{proposition}
\leanref[, which gives (3) and the inclusion $\supseteq$ of (1); the inclusion $\subseteq$ is \path{FurioLombardo.M4.four_dvd_imv_of_mem} in \path{FurioLombardo/M4/Main.lean}, (2) is \path{FurioLombardo.Discharge.M3a.quarter_of_cert} in \path{FurioLombardo/Discharge/M3a/Quarter.lean}, and the data checks are \path{FurioLombardo.M4.T0.checks} and \path{FurioLombardo.M4.T1.checks} in \path{FurioLombardo/M4/Instances.lean}]{FurioLombardo.M4.qchar_of_cert}{FurioLombardo/M4/Cover.lean}

\begin{proof}
The values $\lambda_k(D_i)$, $a$ and $b$ are enclosed in balls around integer vectors $\ell_i$, $\ell_a$, $\ell_b$. For $D_i$, a chain of doublings and additions, computed with the composition and reduction of Cantor \cite{Cantor87} in $2$-adic ball arithmetic, encloses the class $2^JnD_i+E_0$ for an integer $n$ and a large $J$, which lies deep in the chart, and a proved bound on the formal logarithm gives $\lambda_k(D_i)$ modulo $2^3$ at least; $a$ and $b$ are computed in the same way, to precision $2^{47}$ at least for $k=0$ and $2^{52}$ at least for $k=1$. The rest is a finite check on integer matrices. An integer $6\times7$ matrix $C$ with $\sum_iC_{ji}\ell_i\equiv4e_j\pmod8$ gives $4\Zt^6\subseteq\Lambda_k$ by Nakayama's lemma. Integer matrices $G_k$, $H_k$ with $H_kG_k=4I$, $G_k\ell_i\equiv0\pmod4$, and columns of $H_k$ congruent modulo $4$ to combinations of the $\ell_i$, give the description of $\Lambda_k$ in~(1). Let $U$ be an integer matrix such that the last four coordinates of $UG_k\ell_a$ and $UG_k\ell_b$ are divisible by~$2^6$, and let $A_0,A_1$ and $B_0,B_1$ be their first two coordinates. The certificate gives $A_0=4A_0'$ with $A_0'$ odd, $B_0=4B_0'$, and $A_0B_1-A_1B_0$ of $2$-adic valuation exactly~$6$. For (2), the vector $A_0'b-B_0'a$ then lies in $4\Lambda_k$, which defines $r$ with $4r=A_0'b-B_0'a$; the $2\times2$ minor of $a$ and $r$ in the first two coordinates shows that $\Zt a+\Zt r$ is saturated in $\Lambda_k$, and it contains $b$ because $A_0'$ is a unit, so $S_k\subseteq\Zt a+\Zt r$. For (3), explicit relations $2^{e_j}s_j\equiv\alpha_ja+\beta_jb$ for two columns $s_j$ of an integer matrix, with $(e_1,e_2)=(2,2)$ for $k=0$ and $(0,2)$ for $k=1$, determine $S_k$, and $P_k$ is formed by four rows of~$UG_k$.
\end{proof}

A PARI/GP computation of the logarithms to precision $\pi^{135}$ at least gives more: $[\Zt^6:\Lambda_k]=2^7$, the submodule $\Zt a+\Zt b$ has elementary divisors $1$ and $4$ in $\Lambda_k$, and the image of $W_k$ has dimension $3$ in $\Lambda_k/2\Lambda_k$ and dimension $1$ in $\Lambda_k/(S_k+2\Lambda_k)$ (Section~\ref{sec:comp}). The proof does not use these numbers.

\section{The covering of \texorpdfstring{$C(\Qt)$}{C(Q2)}}\label{sec:cover}

This section checks condition (iii) at every lift of a rational point. It covers $C(\Qt)$ by boxes on which either no point lifts to $D_{\delta_k}$ over $\kv$, or the logarithm has a constant leading level and class, or the point is close to a known lift.

\begin{lemma}\label{lem:discs}
Every point of $C(\Qt)$ has homogeneous coordinates of one of the forms
\[
(2X,2Y,1),\quad(1+2X,2Y,1),\quad(1+2X,1+2Y,1),\quad(2X,1,2Y),\quad(1+2X,1,2Y)
\]
with $X,Y\in\Zt$; we call these the discs $1$ to $5$ and $X$ the parameter of the point.
\end{lemma}
\leanref[; a lift of a rational point lies over a point of a disc by \path{liftTwist_Mmat} in \path{FurioLombardo/Discharge/M4Cert/Bridge.lean}]{FurioLombardo.Discharge.M4Cert.exists_disc}{FurioLombardo/Discharge/M4Cert/Discs.lean}

\begin{proof}
Scale a point to a primitive vector over $\Zt$. Its reduction is a point of $C$ over $\F_2$, and the points $(0:1:1)$ and $(1:0:0)$ of $\mathbb P^2(\F_2)$ are not on the reduction of $C$. The five remaining points of $\mathbb P^2(\F_2)$ give the five forms.
\end{proof}

A point of $D_{\delta_k}(\kv)$ over a point $p$ of $C$ is $(p:r:s)$ with $\delta_kr^2=Q_1(p)$, $\delta_krs=Q_2(p)$ and $\delta_ks^2=Q_3(p)$, determined by $p$ up to the involution $\iota$. The certificate of the twist $k$ lists boxes $\{X\equiv c\pmod{2^s}\}$ in the five discs, each with a reference value that selects one of the two points over each point of the box: the one whose coordinate $r$ (or $s$) is closer to the reference. For $X$ in a box let $x_k(d,X)$ be the selected point over the point of parameter $X$ of the disc $d$, when it exists, and
\[
\ell_k(d,X)=\lambda_k\bigl(\varphi_k(x_k(d,X))-\varphi_k(x_a)\bigr),
\]
with $\varphi_k$ taken over $\kv$. The boxes are of three kinds: excluded, constant and tail. For $k=0$ there are $11$ excluded, $82$ constant and $2$ tail boxes, $95$ in all; for $k=1$ there are $12$, $22$ and $2$, $36$ in all. A finite check shows that in each disc the boxes cover $\Zt$.

\begin{lemma}\label{lem:exclusion}
No point in an excluded box is the image of a point of $D_{\delta_k}(\kv)$.
\end{lemma}
\leanref[ and \path{hExcl_T1_Mmat}]{FurioLombardo.Discharge.M4Cert.hExcl_T0_Mmat}{FurioLombardo/Discharge/M4Cert/Bridge.lean}

\begin{proof}
For every pair of residues $(X,Y)$ modulo $4$ in the box at which the equation of $C$ holds modulo $8$, the residue modulo $8$ of $Q_1(p)/\delta_k$, or of $Q_3(p)/\delta_k$, in $\Ov$ is not the residue of a square; the residues are computed from exact data. A point of $D_{\delta_k}(\kv)$ over $p$ would make $Q_1(p)/\delta_k=r^2$ and $Q_3(p)/\delta_k=s^2$ squares in~$\kv$.
\end{proof}

\begin{lemma}\label{lem:constant-box}
Each constant box $\{c+2^sY:Y\in\Zt\}$ of the twist $k$ carries integers $v_M$, divisible by $3$, and $\nu_0$ with $v_M/3+s=\nu_0+3$, such that $\ell_k(d,c+2^sY)-\ell_k(d,c)\in2^{\nu_0+3}\Zt^6$ for every $Y\in\Zt$. One has $v_M=9$ at every constant box for $k=0$ and $v_M=6$ for $k=1$.
\end{lemma}
\leanref[; the per box certificates are \path{BoxData.constCert_0} and \path{BoxData.constCert_1} in \path{FurioLombardo/Discharge/M4Box/BoxData}]{FurioLombardo.Discharge.M4Box.hConstLip_of_cert}{FurioLombardo/Discharge/M4Box/BoxCert.lean}

\begin{proof}
The difference is $\lambda_k$ of $\varphi_k(x_k(d,c+h))-\varphi_k(x_k(d,c))$ with $h=2^sY$. A computation with Taylor models in $h$, that is, truncated power series with ball coefficients and a bound on the remainder over the box, encloses the point of the disc, the square root that gives the selected lift, the Mumford ideal $\varphi_k(x_k(d,c+h))$, its sum with the fixed class $E_0-\varphi_k(x_k(d,c))$, and the chart coordinates $t(h)$ of the difference. It gives $|t_i(h)|\le|\pi|^D$ and $|(\mathrm At(h))_i|\le|\pi|^{v_M+3s}$ for all $h$, with $D\ge M+4$ and $v_M+3s+M+3\le2D$. The formal logarithm differs from the identity by terms of degree at least~$2$, so $|\lambda_k-\mathrm At|\le|\pi|^{2D-M-3}$ on these chart points, and the two bounds give $\lambda_k\in\pi^{v_M+3s}\Ov^2=2^{v_M/3+s}\Zt^6$.
\end{proof}

\begin{lemma}\label{lem:centres}
Let $c$ be the centre of a constant box, with $\nu_0$ as in Lemma~\ref{lem:constant-box}. Then $\ell_k(d,c)=s_0+2^{\nu_0}z_0+2^{\nu_0+1}l$ as in Lemma~\ref{lem:leading-class}, with $S=S_k$ and $W=W_k$.
\end{lemma}
\leanref[ and \path{CentreData.hCe_1}; the classes are checked by the field \path{centres} of \path{FurioLombardo.M4.TwistChecks}]{FurioLombardo.Discharge.M4Box.CentreData.hCe_0}{FurioLombardo/Discharge/M4Box/CentreData/Values0.lean}

\begin{proof}
The $104$ centres have rational parameters. At each of them the point of the disc, the selected lift, $\varphi_k$ and $\lambda_k$ are computed in ball arithmetic, with every branch decision certified, and give $\ell_k(d,c)$ modulo $2^{q}$ with $q\ge\nu_0+3$. The integer vector $y$ obtained satisfies the divisibility conditions of Proposition~\ref{prop:lattice}(3): $2^{\nu_0+2}$ divides $P_ky$, $2^{\nu_0+3}$ does not, and for each $w\in W_k$, $2^{\nu_0+3}$ does not divide $P_k(y-2^{\nu_0}w)$. This is the claim.
\end{proof}

\begin{lemma}\label{lem:tail-box}
Each tail box is $\{X\equiv X_i\pmod{2^s}\}$ in the disc of a known lift $x_i$ with parameter $X_i$, and $\ell_k(d,X_i)\in S_k$. There are $g\in\Lambda_k$ and integers $s_0\le s$, $v_M$ and $\nu$, with $g$ of the form $s_0'+2^{\nu}z_0+2^{\nu+1}l$ of Lemma~\ref{lem:leading-class} and $\nu+2+s_0\le v_M/3+s$, such that for every $m\ge s$ and every unit $u$ of~$\Zt$,
\[
\ell_k(d,X_i+2^mu)-\ell_k(d,X_i)-2^{m-1}u\,g\in2^{v_M/3+2m-s_0}\Zt^6 .
\]
The values $(s,s_0,v_M,\nu)$ are $(6,2,9,5)$ at $x_0$ and $x_2$, $(6,3,6,3)$ at $x_1$ and $(4,2,6,2)$ at~$x_3$.
\end{lemma}
\leanref[ and \path{BoxData.tailCert_0}, \path{BoxData.tailCert_1}]{FurioLombardo.Discharge.M4Box.hTailQuad_of_cert}{FurioLombardo/Discharge/M4Box/BoxCert.lean}

\begin{proof}
The Taylor model computation of Lemma~\ref{lem:constant-box} over the tail box writes the chart coordinates as $t(h)=c_0+hc_1+h^2R(h)$ with bounds on $c_1$ and on $R$. The constant term $c_0$ vanishes, because at $h=0$ the difference is $0$ and $\psi$ is injective. The vector $g$ is the one with coordinates $2\mathrm Ac_1$ in $\Zt^6$, enclosed in a ball around an integer vector whose coordinates are divisible by $4$, so $g\in4\Zt^6\subseteq\Lambda_k$, and its leading level and class are read with $P_k$ as in Lemma~\ref{lem:centres}. The bound on the logarithm gives the displayed congruence. Finally $\ell_k(d,X_i)$ is $\pm\lambda_k(\varphi_k(x_i))-a$, which lies in~$S_k$.
\end{proof}

\begin{proposition}\label{prop:cover}
Let $P$ be a rational point of $C$, $x\in D_{\delta_k}(K)$ a lift of $P$, and $y=\lambda_k\bigl(\varphi_k(x)-\varphi_k(x_a)\bigr)$, computed in $A_k(\kv)$. Then condition (iii) holds at $y$ with $S=S_k$ and $W=W_k$, and if $y\in S_k$ then $P$ is one of the two known points $\pi(x_a)$, $\pi(x_b)$ of the twist.
\end{proposition}
\leanref[; the step at one parameter is \path{FurioLombardo.Discharge.M4Box.lead_or_known_w} in \path{FurioLombardo/Discharge/M4Box/Chain.lean}, the known points at the parameters $X_i$ are \path{FurioLombardo.Discharge.M4Cert.tailGood_T0_Mmat} and \path{FurioLombardo.Discharge.M4Cert.tailGood_T1_Mmat} in \path{FurioLombardo/Discharge/M4Cert/Bridge.lean}, and the relation between $y$ and $\ell_k$ is \path{hLog_of_cover} in \path{FurioLombardo/Discharge/M4Box/LiftLog.lean}]{FurioLombardo.Discharge.M4Box.coverIII_knownZero_of_M4_w}{FurioLombardo/Discharge/M4Box/ChainM3a.lean}

\begin{proof}
By Lemma~\ref{lem:discs}, $P$ is the point of parameter $X$ of a disc $d$, and $X$ lies in a box, which is not excluded by Lemma~\ref{lem:exclusion}. The image of $x$ in $D_{\delta_k}(\kv)$ is $x_k(d,X)$ or $\iota x_k(d,X)$ up to rescaling, and $\varphi_k(\iota x')=-\varphi_k(x')$, so $y=\ell_k(d,X)$ or $y=-\ell_k(d,X)-2a$; since $a\in S_k$, the form of Lemma~\ref{lem:leading-class} passes from $\ell_k(d,X)$ to $y$. If $X$ is in a constant box with centre $c$, Lemmas~\ref{lem:constant-box} and~\ref{lem:centres} give $\ell_k(d,X)=s_0+2^{\nu_0}z_0+2^{\nu_0+1}l+e$ with $e\in2^{\nu_0+3}\Zt^6\subseteq2^{\nu_0+1}\Lambda_k$ by Proposition~\ref{prop:lattice}(1), so $\ell_k(d,X)$ has the form of Lemma~\ref{lem:leading-class}, which gives condition (iii) and $y\notin S_k$. If $X$ is in a tail box and $X\neq X_i$, write $X=X_i+2^mu$ with $m\ge s$ and $u$ a unit; by Lemma~\ref{lem:tail-box}, $\ell_k(d,X)$ is an element of $S_k$ plus $2^{m-1}ug$ plus an element of $2^{v_M/3+2m-s_0}\Zt^6$, which lies in $2^{\nu+m}\Lambda_k$ because $v_M/3+2m-s_0-2\ge\nu+m$, and $2^{m-1}ug$ has leading level $\nu+m-1$ and the leading class of $g$; the conclusion is the same. If $X=X_i$, the equation of $C$ at the parameter $X_i$ has the single root $Y=0$, so $P$ is the image of the known lift~$x_i$.
\end{proof}

\section{Proof of Theorem~\ref{thm:main}}\label{sec:proof}

\begin{theorem}\label{thm:twist}
For $k=0,1$, every rational point of $C$ that lifts to $D_{\delta_k}(K)$ is $P_0$ or $P_2$ if $k=0$, and $P_1$ or $P_3$ if $k=1$.
\end{theorem}
\leanref[; it is applied to each twist in \path{FurioLombardo/Discharge/M4Box/ChainM3a.lean}, and the twist $k$ uses Theorem~\ref{thm:selmer-basis} for that twist]{FurioLombardo.M3a.Route.twistConclusion_of_inputs}{FurioLombardo/M3a/Route.lean}

\begin{proof}
Apply Corollary~\ref{cor:stoll-sat} with $\mathsf A=A_k(K)$, $\mathsf B=A_k(\kv)$, $\iota$ the map induced by $K\to\kv$, $\lambda=\lambda_k$, $T=T_k$, $\gamma_a=\varphi_k(x_a)$, $\gamma_b=\varphi_k(x_b)$, $S=S_k$ and $W=W_k$, to $\varphi_k(x)-\varphi_k(x_a)$ for a lift $x$ of a rational point~$P$. Condition (loc) is Corollary~\ref{cor:selmer-consequences}(1), and (ker) is Lemma~\ref{lem:local-divisors}. The condition on $W$ follows from Corollary~\ref{cor:selmer-consequences}(2), which gives $\lambda\iota(Q)=\sum_i(B\epsilon)_i\lambda_k(D_i)+2\lambda_k(b)$, since $\sum_i(B\epsilon)_i(\lambda_k(D_i)-\ell_i)$ lies in $8\Zt^6$, which is contained in $2\Lambda_k$ by Proposition~\ref{prop:lattice}(1). The lattice condition is Proposition~\ref{prop:lattice}(2), and condition (iii) is Proposition~\ref{prop:cover}. So $\lambda_k(\varphi_k(x)-\varphi_k(x_a))\in S_k$, and Proposition~\ref{prop:cover} shows that $P$ is one of the two known points of the twist.
\end{proof}

\begin{proof}[Proof of Theorem~\ref{thm:main}]
The four points lie on $C$ by substitution. Let $P$ be a rational point of $C$. By Corollary~\ref{cor:lift}, $P$ lifts to $D_{\delta_0}(K)$ or to $D_{\delta_1}(K)$, and Theorem~\ref{thm:twist} shows that $P$ is one of $P_0,P_1,P_2,P_3$. Two nonzero vectors give the same point exactly when they are proportional, which is the statement of Theorem~\ref{thm:main}.
\end{proof}

\section{Computations and certificates}\label{sec:comp}

This section says what the Lean proof checks and how, which programs wrote the data it checks, which computations a second program redid, and which computations came before the formal proof. Every script named below is in the repository, every result quoted from a run is in the recorded output of that run, and paths are relative to the root of the repository.

\subsection{The formal proof}\label{sec:comp-lean}
The statement is fixed in the file \path{FurioLombardo/Statement.lean}, whose definitions are the following.
\lstinputlisting[linerange={22-45}]{\statementfile}
The theorem \path{FurioLombardo.conjecture_1_6} in \path{FurioLombardo/Main.lean} has the type \path{FurioLombardo.Conjecture}. It combines \path{FurioLombardo.Discharge.M4Box.onlyFourPoints} in \path{FurioLombardo/Discharge/M4Box/Final.lean}, of type \path{FurioLombardo.OnlyFourPoints}, with the equivalence \path{FurioLombardo.conjecture_iff_onlyFourPoints} of the statement file. The development uses Lean~4 \cite{Lean4}, version 4.34.1, and Mathlib \cite{Mathlib20} at its tag v4.34.1. The command \texttt{\#print axioms} applied to \path{FurioLombardo.conjecture_1_6} prints \texttt{propext}, \texttt{Classical.choice} and \texttt{Quot.sound}, the standard axioms of Lean and Mathlib, and nothing else. The continuous integration of the repository builds the development from source, runs this command on \path{FurioLombardo.conjecture_1_6} and on \path{FurioLombardo.Discharge.M4Box.onlyFourPoints}, and fails if any other axiom appears. It also runs Comparator \cite{Comparator}, which replays the proof in the Lean kernel and checks that \path{FurioLombardo.conjecture_1_6} proves the statement of \path{FurioLombardo/Challenge.lean}, a file that imports only Mathlib and repeats the definitions of the statement file, with no axiom other than these three. The files of a release are attached only after these checks pass.

Each finite verification in Sections~\ref{sec:descent} to~\ref{sec:cover} is proved in Lean by \texttt{decide +kernel}, \texttt{decide}, \texttt{norm\_num} or \texttt{rfl}, so the Lean kernel checks it. The typical form is the equality of a Boolean function of explicit data with \texttt{true}, proved by \texttt{decide +kernel}, which has the kernel evaluate the function. No proof of the development uses \texttt{native\_decide}, so no compiled code belongs to the trusted base; the continuous integration checks the sources for it, and in this version of Lean each use of it would add an axiom of its own to the list above (\path{code/formal-proof/native_decide_control.out}). For each such function, a lemma proved in Lean states what the value \texttt{true} means: that an ideal is principal, that a square class is nontrivial, that a vector of $\Zt^6$ lies in the lattice $\Lambda_k$. The $2$-adic quantities of Sections~\ref{sec:local} and~\ref{sec:cover} are computed with a ball arithmetic over $\kv$ formalized in Lean. A ball is given by three integers $c_0,c_1,c_2$, a scale $e$, a radius $r$ and a lower bound for the valuation of its centre; it contains $x\in\kv$ when $|2^ex-(c_0+c_1\pi+c_2\pi^2)|\le|\pi|^r$. Each operation comes with a proof that the ball it returns contains the result of the operation on elements of the input balls, square roots are certified by Hensel's lemma, and the Taylor models of Lemma~\ref{lem:constant-box} are polynomials with ball coefficients and a ball for the remainder.

Besides Mathlib, the development uses code taken from four Lean projects, all under the Apache License~2.0: from Michael Stoll's EllipticCurves \cite{StollEC}, formal group laws and their logarithms, fractional ideals, $S$-integers and Selmer groups; from Tau Ceti \cite{TauCeti}, results on local fields (uniformizers, normalized valuations, squares) and the integral closedness of the coordinate ring, adapted from the case of Weierstrass equations to $Y^2=f$; from Chris Birkbeck's AINTLIB \cite{AINTLIB}, the completed Dedekind zeta function; and from Lean Pool \cite{LeanPool}, contour integrals over rectangles, in a file of the FLT Project adapted from the PNT+ project of Kontorovich and Tao. The last two enter the proof of Zimmert's bound in Proposition~\ref{prop:classnumber}. The repository names, for each file taken, its source, the commit or dated snapshot it was taken from, and the changes made.

\subsection{The data of the checks}\label{sec:comp-data}
The data that the kernel checks were written by PARI/GP scripts and, for the values of $\lambda_k$ at the points $D_i$ and at the known lifts, the boxes and the centres, by Lean programs run outside the proof. The proof does not trust these programs, since the kernel checks what they write. Table~\ref{tab:data} lists them by section. The script \path{code/field/class_number_primes.gp} writes the $11298$ certificates of the proof of Proposition~\ref{prop:classnumber}, one for each prime $p\le120000$ other than $2$, $7$ and $45613$ (\path{code/field/class_number_primes_rerun.out}).
\begin{table}[ht]
\centering\small
\begin{tabular}{>{\raggedright\arraybackslash}p{0.34\textwidth} >{\raggedright\arraybackslash}p{0.58\textwidth}}
\toprule
Part of the proof & Programs, in \path{code/}\\
\midrule
Section~\ref{sec:descent} & \path{field}, \path{descent}\\
Section~\ref{sec:prym} & \path{genus2-curves}\\
Section~\ref{sec:selmer}, global bound & \path{selmer-global-bound}\\
Lemma~\ref{lem:local-images} & \path{selmer-local-conditions} (the places $v$, $w_{12}$, $w_6$, $w_7$ and the real places)\\
Theorem~\ref{thm:selmer-basis} & \path{selmer-assembly}\\
Section~\ref{sec:local} & \path{local-group}, the lattice programs of \path{covering}, and \path{local-group/di_value_certs_gen.lean} and \path{local-group/known_lift_certs_gen.lean} (Lean)\\
Section~\ref{sec:cover} & \path{covering}, with \path{covering/box_certs_gen.lean} and \path{covering/centre_certs_gen.lean} (Lean)\\
\bottomrule
\end{tabular}
\caption{Programs that wrote the data of the Lean checks, by directory of \texttt{code/}; ``Lean'' marks the Lean programs.}\label{tab:data}
\end{table}
Several of these programs start from exact data written by PARI/GP scripts of \path{code/earlier-computations}: the field $K$, the forms $Q_i$ and the polynomials $f_\delta$ (\path{bruin_form.gp}), the field $\kv$ (\path{completion_kv.gp}), the discs of Lemma~\ref{lem:discs} (\path{discs_q2.gp}), the values of $\varphi_k$ at the known lifts (\path{phi_known_lifts.gp}), and the lattice data, the boxes and the centres of Table~\ref{tab:runs}.

\subsection{Earlier computations}\label{sec:comp-other}
The data of the proof were first computed and checked in PARI/GP, before the formal proof existed. Table~\ref{tab:runs} lists these computations with their running times; the Lean proof checks their results again.
\begin{table}[ht]
\centering\small
\begin{tabular}{>{\raggedright\arraybackslash}p{0.38\textwidth} >{\raggedright\arraybackslash}p{0.36\textwidth} r}
\toprule
Computation & Script in \path{code/earlier-computations} & Time\\
\midrule
$\Cl(\OK)$ trivial, all prime ideals up to the Minkowski bound & \path{class_number_minkowski.gp}, $8$ ranges, $4$ at a time & $8194$~s\\
Boxes of Section~\ref{sec:cover} & \path{boxes.gp} & $539+151$~s\\
Lattice data of Proposition~\ref{prop:lattice} and of the remark after it & \path{lattice_cert.gp} & $4+4$~s\\
Leading level and class at the centres (Lemma~\ref{lem:centres}) & \path{centres_cert.gp} & $27+9$~s\\
\bottomrule
\end{tabular}
\caption{PARI/GP computations. Two times are those of the twists $k=0$ and $k=1$; the first time is the sum of the eight ranges.}\label{tab:runs}
\end{table}

The first line of the table is the computation mentioned after Proposition~\ref{prop:classnumber}: each of the eight ranges of rational primes up to the Minkowski bound took between $957$~s and $1153$~s, and every prime ideal of norm at most that bound was shown principal by an explicit generator. The proof uses Zimmert's bound instead.

The script \path{code/selmer-local-conditions/selmer_rows.gp} computes, for a twist $k$ and a place $w$, the matrix $C_w$ of Lemma~\ref{lem:local-images}, and prints a check at each step. It factors $f_k$ over $K_w$ into the factors listed in the proof of the lemma, with a Hensel certificate for a root of each factor, tests its square class coordinates on known answers and on negative controls, computes the coordinates of generators of $K_w^\times$ modulo squares, of the images under $\mu_w$ of $n_w$ explicit points and of the $82$ classes $g_s$, checks that the $n_w$ images are independent modulo the image of $K_w^\times$, and writes the rows of $C_w$. Its outputs at $w_{12}$ (\path{selmer_rows_twist0_w12.out} and \path{selmer_rows_twist1_w12.out}) and at $w_7$ (\path{selmer_rows_twist1_w7.out}), in the same directory, report no failed check; these runs took $172$~s, $172$~s and $10$~s. A last stage, \path{selmer_rows_kernel.out}, stacks the rows of all the places of $\Sigma_k$ and finds a kernel of dimension $19$ for each twist, as in the proof of Theorem~\ref{thm:selmer-basis}; it takes one second. At $w_{12}$ and $w_7$ the Lean proof checks that the linear forms certified from the data of \path{code/selmer-local-conditions/local_data_w12.gp} and \path{code/selmer-local-conditions/local_data_w7.gp} give exactly these rows.

The route of Remark~\ref{rem:alt-descent} was checked by computation only: the $2$-descent is in \path{code/jacobian-over-q} (the span of the fake Selmer group, the kernel of the descent map and the rank in \path{known_classes_images.out}, with the class group and units of the number field of the descent certified in \path{fake_selmer_group.out}), the orders of $J(\F_3)$ and $J(\F_{11})$ in \path{code/jacobian-over-q/jacobian_order_f3_f11.out}, and the Mordell-Weil sieve, with $233$ primes up to $1999$, in \path{code/jacobian-over-q/mw-sieve/mwsieve_p2000_all_firstfit.txt}.

\subsection{Second implementations}\label{sec:comp-second}
The certificates of Sections~\ref{sec:local} and~\ref{sec:cover} were written by one set of programs. To guard against a bug in them, a second program, written separately in Sage \cite{Sage} from the mathematics alone, redid that part in \path{code/second-implementations/local-group-covering}. It takes from text exports of the first programs the exact global data, the points $D_i$, the list of boxes, the image $W_k$ of the Selmer group at $v$ and the matrix $\mathrm A$ of Proposition~\ref{prop:log}, and recomputes the logarithms of the $D_i$ and of the known points to precision $\pi^{1486}$ at least, the lattice $\Lambda_k$ and the saturation $S_k$, the leading level and class at each of the $104$ centres, a bound $v_M$ on each box by its own Taylor bounds, the covering of each disc and the $23$ exclusions. It agrees with the first programs at every centre and every box. It does not redo the independence of the $\mu_v(D_i)$ (Lemma~\ref{lem:local-divisors}) or Lemma~\ref{lem:local-2-torsion}. The box data of the Lean proof are written by \path{code/covering/lattice_boxes_data.gp} from the same box lists; the lattice matrices of the Lean data are not compared with this program. Its logarithms computed at precisions $400$ and $1000$ agree to the precision claimed at $400$, and the centres took $22.5$~s at precision $400$ and $216.6$~s at precision~$1000$.

A Sage program written separately, in \path{code/second-implementations/descent-set}, recomputes the descent set of Theorem~\ref{thm:selmer-set}: the group $K(S,2)$, the sieve of Lemma~\ref{lem:sieve}, and the local conditions at $2$ and at $3$ by its own coverings of $C(\Qt)$ and $C(\Q_3)$, in which the descent class is constant on each of $40$ and $20$ boxes and the other boxes contain no point of~$C$. It takes the class number $h_K=1$ as an input, finds its $S$-units under the generalized Riemann hypothesis and then checks them without it. It keeps the same two classes $\delta_0$ and $\delta_1$, and the $1120$ local square decisions of its coverings agree with those of the PARI/GP function \texttt{nfislocalpower}.

For the Selmer bound of Section~\ref{sec:selmer}, an earlier PARI/GP computation shares its global data with \path{selmer_rows.gp} but computes the local images with other code (\path{code/earlier-computations/prym_two_descent.gp}); it also finds, for both twists, a space of dimension $19$ of vectors $a\in\F_2^{82}$ that satisfy all the local conditions (183~s, starting from cached intermediate data). The Sage script \path{code/second-implementations/selmer-space/selmer_space_check.sage} recomputes that space over $\F_2$ from the generator coordinates and the local images exported by this earlier computation, and finds the same dimension and the same image at~$v$; it takes the local images as given and checks the linear algebra.

\subsection{Software, hardware and data}
The computations used PARI/GP 2.17.4 \cite{PARI}, Sage 10.9 \cite{Sage}, and Lean~4.34.1 with Mathlib~v4.34.1, on one machine with an AMD Ryzen~9 5900X processor ($12$ cores, $24$ threads) and $62$~GB of memory (\path{paper/machine.txt}). The repository \url{\repo} contains the Lean development, the PARI/GP and Sage scripts with their complete outputs, the certificates, and the \TeX{} source of this paper.

\section*{Contact}
Questions, corrections and comments are welcome as issues on the repository: \url{\repo/issues/new/choose}.

\bibliographystyle{abbrv}
\bibliography{refs}

\end{document}
